\chapter{Supervised Learning Core}
\label{cml:chap:supervised_core}

\begin{tldr}
This is the derivation backbone of Volume IV. Linear and logistic regression are the classical loop's ``can you do mathematics on a whiteboard'' probe, and the regularization-geometry question is the single most-asked classical question after trees. The chapter builds, in order: linear regression from the normal equations, with the four reasons the closed form fails and why no serious library ever forms $(X^{\top}X)^{-1}$; ridge, lasso, and elastic net with \emph{both} halves of the sparsity answer---the constraint-set geometry picture and the soft-thresholding algebra that turns ``L1 gives sparsity'' from a slogan into a theorem; logistic regression from the Bernoulli likelihood to the famously clean gradient $X^{\top}(\sigma(X\beta)-y)$, the convexity and IRLS story, separation as the classic blow-up, and the GLM frame that explains why that gradient shape is not a coincidence; the SVM from geometric margin through the Lagrangian dual to the kernel trick, with support vectors read off the KKT conditions and an honest note on why SVMs are still \emph{asked} in 2026 despite rarely being \emph{shipped}; Naive Bayes as the generative half of the canonical generative--discriminative pair; and $k$-NN as the cleanest place to watch the curse of dimensionality bite. The bias--variance decomposition itself is canonical in [DL 2]---here it is instantiated as a named knob on every model. Tree ensembles, the other half of the classical canon, are Chapter~\ref{cml:chap:trees_ensembles}.
\end{tldr}

%==============================================================================
\section{Linear Regression}
\label{cml:sec:linear_regression}
%==============================================================================

Linear regression is the most over-asked and under-prepared topic in the classical loop. Everyone can write $y = X\beta + \varepsilon$; far fewer can derive the normal equations without hesitation, state which assumption buys which property, name four distinct reasons the closed form is the wrong implementation, or explain what a library actually does instead. That gap is exactly what the question is for.

\subsection{The Model and What Each Assumption Buys}

The model is $y_i = x_i^{\top}\beta + \varepsilon_i$, or in matrix form $y = X\beta + \varepsilon$ with $X \in \R^{n \times d}$ (a column of ones absorbs the intercept). Ordinary least squares picks $\beta$ to minimize the residual sum of squares $\norm{y - X\beta}_2^2$.

The assumptions are worth separating by \emph{what they are for}, because interviewers probe exactly this seam:

\begin{itemize}
    \item \textbf{Linearity in the parameters.} The model must be linear in $\beta$, not in $x$: $y = \beta_0 + \beta_1 x + \beta_2 x^2$ is a linear model. Violating this biases everything.
    \item \textbf{Exogeneity}, $\E[\varepsilon \given X] = 0$. This is the one that buys \emph{unbiasedness}. It fails under omitted confounders, simultaneity, and measurement error in $X$---the entire subject of Chapter~\ref{cml:chap:experimentation_causal}.
    \item \textbf{Homoscedasticity and no autocorrelation}, $\Var(\varepsilon \given X) = \sigma^2 I$. These buy \emph{efficiency}, not unbiasedness: violate them and OLS still points at the right answer, but it is no longer the minimum-variance linear estimator and the usual standard errors are wrong. Fixes: robust (Huber--White) standard errors, clustered errors, or GLS/weighted least squares.
    \item \textbf{Full column rank} of $X$. This buys \emph{identifiability}---without it $\beta$ is not unique.
    \item \textbf{Normality of the errors}, $\varepsilon \sim \mathcal{N}(0, \sigma^2 I)$. This buys \emph{exact finite-sample inference} ($t$ and $F$ distributions) and makes OLS coincide with the Gaussian MLE. It is not needed for OLS to be unbiased, consistent, or BLUE, and asymptotically the CLT supplies the inference anyway.
\end{itemize}

\begin{keyinsight}[Gauss--Markov: Normality Is Not on the List]
Under linearity, exogeneity, homoscedasticity, and no autocorrelation, OLS is the \textbf{Best Linear Unbiased Estimator}---minimum variance among all linear unbiased estimators. Normality appears nowhere in that statement. Two consequences a candidate should have ready: (i) ``OLS requires normally distributed data'' is wrong twice over---normality is about the \emph{errors}, never the features, and it is about \emph{inference}, not estimation; (ii) ``best'' is quantified over \emph{linear unbiased} estimators only, which is exactly why ridge---a \emph{biased} estimator, and therefore simply outside the comparison class---can beat OLS on mean squared error (Section~\ref{cml:sec:regularization}).
\end{keyinsight}

\subsection{Deriving the Normal Equations}

This is the single most requested classical derivation. It should take ninety seconds and no notes.

\begin{mathresult}{The Normal Equations}
Write the objective and expand:
\begin{equation}
J(\beta) = \norm{y - X\beta}_2^2 = (y - X\beta)^{\top}(y - X\beta) = y^{\top}y - 2\beta^{\top}X^{\top}y + \beta^{\top}X^{\top}X\beta .
\end{equation}
Differentiate, using $\nabla_\beta (\beta^{\top}A\beta) = 2A\beta$ for symmetric $A$ and $\nabla_\beta (b^{\top}\beta) = b$:
\begin{equation}
\nabla_\beta J = -2X^{\top}y + 2X^{\top}X\beta .
\end{equation}
Set to zero to obtain the \textbf{normal equations}, and solve when $X^{\top}X$ is invertible:
\begin{equation}
X^{\top}X\hat\beta = X^{\top}y \qquad\Longrightarrow\qquad \hat\beta = (X^{\top}X)^{-1}X^{\top}y .
\end{equation}
The Hessian is $\nabla^2_\beta J = 2X^{\top}X \succeq 0$, so $J$ is convex and the stationary point is a global minimum; it is \emph{strictly} convex, hence the minimizer is unique, exactly when $X$ has full column rank.
\end{mathresult}

Two things to say out loud while writing it. First, the name: the equations assert $X^{\top}(y - X\hat\beta) = 0$, i.e., the residual is \textbf{orthogonal} (normal) to every column of $X$. Second, the same first-order condition with a general convex loss is what produces the whole GLM family in Section~\ref{cml:sec:logistic_glm}; the squared-loss case is special only in that it can be solved in closed form.

\subsection{The Geometry: Projection, the Hat Matrix, and Leverage}

The geometric statement is often the fastest way to answer follow-ups. Fitted values are
\begin{equation}
\hat y = X\hat\beta = \underbrace{X(X^{\top}X)^{-1}X^{\top}}_{H,\ \text{the ``hat matrix''}} y ,
\end{equation}
and $H$ is the \textbf{orthogonal projection} onto the column space of $X$: symmetric, idempotent ($H^2 = H$), with $\rank(H) = \tr(H) = d$. OLS therefore does exactly one thing---drop a perpendicular from $y$ onto the subspace spanned by the features---and the residual $y - \hat y = (I-H)y$ lives in the orthogonal complement. This is why the normal equations are an orthogonality condition, and why adding a feature can never increase training RSS (you are projecting onto a larger subspace).

The diagonal $h_{ii}$ is observation $i$'s \textbf{leverage}: how much $\hat y_i$ moves when $y_i$ moves, since $\partial \hat y_i / \partial y_i = h_{ii}$. Useful facts: $0 \le h_{ii} \le 1$, $\sum_i h_{ii} = d$, so the average leverage is $d/n$ and the standard screen is $h_{ii} > 2d/n$. Leverage is a property of $X$ alone---an outlier in feature space---while \emph{influence} (Cook's distance) combines leverage with a large residual. And the leave-one-out residual has the closed form $y_i - \hat y_i^{(-i)} = (y_i - \hat y_i)/(1 - h_{ii})$, which is why LOOCV for linear models costs one fit rather than $n$.

\subsection{Closed Form Versus Gradient Descent}

\begin{mathresult}{Cost of the Two Routes}
\textbf{Normal equations.} Forming $X^{\top}X$ costs $O(nd^2)$; factorizing and solving costs $O(d^3)$ (Cholesky is $\tfrac{1}{3}d^3$; QR on $X$ directly is $\approx 2nd^2 - \tfrac{2}{3}d^3$). Total $O(nd^2 + d^3)$, and memory $O(d^2)$ for the Gram matrix.

\textbf{Gradient descent.} Each step costs one matrix--vector product each way, $O(nd)$; with $k$ steps, $O(ndk)$. Memory is $O(d)$ if you stream $X$, and stochastic/mini-batch variants cut the per-step cost further.

\textbf{The crossover} is $d^2 + d^3/n$ versus $dk$: the closed form wins when $d$ is small relative to the number of iterations you would need, and loses badly once $d$ is large.
\end{mathresult}

Worked, because interviewers want the arithmetic. At $n = 10^5$, $d = 10^3$: $nd^2 = 10^{11}$ and $d^3 = 10^9$, a second or two---take the closed form. At $n = 10^6$, $d = 10^4$: $nd^2 = 10^{14}$ flops, and $X^{\top}X$ alone is $10^8$ doubles, $800$\,MB. Gradient descent at $k=100$ passes is $ndk = 10^{12}$ flops, two orders of magnitude cheaper, with $O(d)$ memory. The closed form is not merely slower here; it does not fit.

\begin{warningbox}
\textbf{Four distinct reasons the closed form fails}---name them separately, because ``it is slow'' is only one of them:
\begin{enumerate}
    \item \textbf{Cost/memory}: $O(nd^2 + d^3)$ time and $O(d^2)$ memory become prohibitive for large $d$.
    \item \textbf{Singularity}: if $d > n$, or any feature is an exact linear combination of others (a one-hot encoding with all levels plus an intercept, a duplicated column), $X^{\top}X$ is singular and $(X^{\top}X)^{-1}$ does not exist. The minimizer is then a non-unique affine set.
    \item \textbf{Ill-conditioning}: near-collinearity makes $X^{\top}X$ nearly singular. Note $\kappa(X^{\top}X) = \kappa(X)^2$: forming the Gram matrix \emph{squares} the condition number, so you lose roughly half your significant digits before you have solved anything.
    \item \textbf{The data does not sit still}: streaming or online settings, data too large for one machine, or a loss that is not squared error at all (logistic, hinge, quantile) admit no closed form regardless.
\end{enumerate}
\end{warningbox}

\subsection{What Libraries Actually Do}

``Compute \texttt{inv(X.T @ X) @ X.T @ y}'' is a red-flag answer even when it runs. Production solvers never form the inverse:

\begin{itemize}
    \item \textbf{Cholesky} on the normal equations ($X^{\top}X = LL^{\top}$, then two triangular solves). Fastest when $d \ll n$ and the problem is well conditioned; inherits the squared condition number.
    \item \textbf{QR} on $X$ itself ($X = QR$, solve $R\hat\beta = Q^{\top}y$). Never forms $X^{\top}X$, so it works at $\kappa(X)$ rather than $\kappa(X)^2$. This is the default for well-posed problems in R and most numerical libraries.
    \item \textbf{SVD / pseudoinverse} ($X = UDV^{\top}$, $\hat\beta = V D^{+} U^{\top} y$). The most robust and the most expensive. It is defined even when $X$ is rank-deficient, where it returns the \emph{minimum-norm} least-squares solution---which is why the \texttt{lstsq}-style solvers in NumPy and scikit-learn silently return an answer for $d > n$ instead of raising. Small singular values can be truncated (\texttt{rcond}), which is a crude cousin of ridge.
\end{itemize}

The one-sentence version for an interview: \textit{solve the least-squares problem, do not invert a matrix; use QR by default and SVD when you suspect rank deficiency, because forming $X^{\top}X$ squares the conditioning.}

\subsection{Multicollinearity}

Collinearity does not bias OLS and does not hurt prediction much within the training distribution. What it destroys is the \emph{variance} of individual coefficients, and therefore every interpretation you wanted to draw from them.

\begin{mathresult}{Variance Inflation}
For predictor $j$, let $R_j^2$ be the $R^2$ from regressing $x_j$ on all the other predictors. Then
\begin{equation}
\Var(\hat\beta_j) = \frac{\sigma^2}{(n-1)\,s_j^2}\cdot \underbrace{\frac{1}{1 - R_j^2}}_{\text{VIF}_j},
\end{equation}
where $s_j^2$ is the sample variance of $x_j$. The variance inflation factor multiplies the variance you would have had if $x_j$ were orthogonal to everything else.
\end{mathresult}

The arithmetic makes it vivid. Two predictors correlated at $\rho = 0.9$ give $\text{VIF} = 1/(1-0.81) = 5.3$, a $2.3\times$ inflation of the standard error. At $\rho = 0.99$, $\text{VIF} = 1/(1 - 0.9801) = 50.3$, a $7.1\times$ inflation---the coefficient can flip sign between bootstrap resamples while predictions barely move. Conventional flags are $\text{VIF} > 5$ (watch) and $> 10$ (act); the condition number of the scaled design, $\kappa(X)$, is the global counterpart, with $> 30$ the usual warning line.

Remedies, in the order they should be offered:
\begin{enumerate}
    \item Decide whether you actually need the coefficients. If the goal is pure prediction, collinearity is largely a non-issue.
    \item Drop or combine the redundant features---a sum, a ratio, a principal component.
    \item Use \textbf{ridge}, the principled fix, because it directly repairs the conditioning (Section~\ref{cml:sec:regularization}).
    \item Get more data, which shrinks the $\sigma^2/[(n-1)s_j^2]$ factor but never touches the VIF itself.
\end{enumerate}

\begin{warningbox}
\textbf{Perfect} collinearity is a different failure from \textbf{near} collinearity. The dummy-variable trap---one-hot encoding a $k$-level categorical into $k$ columns \emph{and} keeping an intercept---makes $X^{\top}X$ exactly singular. Drop one level (or the intercept), or use a regularized fit, which tolerates the redundancy by construction. This is also why penalized models are usually fed the full one-hot encoding without complaint while OLS is not.
\end{warningbox}

\subsection{Diagnostics Worth Naming}

A candidate who can only fit a linear model looks junior next to one who can interrogate it. The short list, with what each is for:

\begin{itemize}
    \item \textbf{Residuals versus fitted values.} Curvature means a missing nonlinearity or interaction; a funnel shape means heteroscedasticity (fix the standard errors, or model the variance, or transform $y$); structure in residuals ordered by time means autocorrelation, which is Chapter~\ref{cml:chap:time_series} territory.
    \item \textbf{Q--Q plot of residuals.} Checks the normality assumption, which matters for small-sample inference and not at all for the point estimate. Heavy tails argue for robust regression (Huber loss) or quantile regression.
    \item \textbf{Leverage and Cook's distance.} Leverage $h_{ii}$ finds unusual $x$; Cook's distance combines it with the residual to find points whose deletion actually moves the fit. Investigate them; do not delete them reflexively.
    \item \textbf{$R^2$ and its traps.} $R^2$ never decreases when a feature is added, so it cannot be used for model selection---use adjusted $R^2$, an information criterion, or, better, held-out error. A high $R^2$ says nothing about causal validity, and a low one can still be a useful model when the irreducible noise is large.
\end{itemize}

%==============================================================================
\section{Regularization: Ridge, Lasso, and Elastic Net}
\label{cml:sec:regularization}
%==============================================================================

\subsection{Ridge Regression}

Ridge adds a squared-$\ell_2$ penalty to the least-squares objective:
\begin{equation}
\hat\beta^{\text{ridge}} = \argmin_\beta \ \tfrac{1}{2}\norm{y - X\beta}_2^2 + \tfrac{\lambda}{2}\norm{\beta}_2^2
\qquad\Longrightarrow\qquad
\hat\beta^{\text{ridge}} = (X^{\top}X + \lambda I)^{-1}X^{\top}y .
\end{equation}
The derivation is one line from the OLS gradient: $-X^{\top}(y - X\beta) + \lambda\beta = 0$.

\begin{keyinsight}[Why Ridge Always Has a Unique Solution]
$X^{\top}X$ is positive semi-definite, so its eigenvalues are $\ge 0$. Adding $\lambda I$ shifts every eigenvalue up by $\lambda$, making them all $\ge \lambda > 0$: $X^{\top}X + \lambda I$ is positive definite, hence invertible, \emph{for any $X$}---including $d \gg n$, exactly duplicated columns, and the dummy-variable trap. Equivalently, the objective is $\lambda$-strongly convex, so the minimizer exists and is unique. This is the cleanest one-sentence answer to ``why does ridge work when $n < d$ and OLS does not,'' and it is also why ridge is the principled fix for multicollinearity: it does not merely shrink coefficients, it \textbf{repairs the conditioning of the problem}, moving $\kappa$ from $d_{\max}^2/d_{\min}^2$ to $(d_{\max}^2+\lambda)/(d_{\min}^2+\lambda)$.
\end{keyinsight}

The SVD view says precisely what ridge does. With $X = UDV^{\top}$ and singular values $d_1 \ge \cdots \ge d_p > 0$,
\begin{equation}
X\hat\beta^{\text{ridge}} = \sum_{j} u_j \frac{d_j^2}{d_j^2 + \lambda} u_j^{\top} y ,
\qquad
\mathrm{df}(\lambda) = \sum_j \frac{d_j^2}{d_j^2+\lambda} .
\end{equation}
Ridge is a \emph{soft filter} on the principal components of $X$: directions with large variance ($d_j^2 \gg \lambda$) pass through essentially untouched, while low-variance directions---the ones where the data barely pins the coefficient down and the estimate is mostly noise---are shrunk toward zero. The effective degrees of freedom $\mathrm{df}(\lambda)$ slides continuously from $p$ at $\lambda=0$ to $0$ as $\lambda\to\infty$, which is the honest way to say ``ridge reduces model complexity'' with a number attached.

\begin{mathresult}{Why a Biased Estimator Beats the BLUE}
Gauss--Markov says OLS is best among \emph{unbiased} linear estimators. Ridge escapes by giving up unbiasedness, and the one-dimensional case shows exactly what it buys. With $X^{\top}X = I$ and a single coefficient, $\hat\beta^{\text{OLS}} \sim \mathcal{N}(\beta, \sigma^2)$ and $\hat\beta^{\text{ridge}} = \hat\beta^{\text{OLS}}/(1+\lambda)$, so
\begin{equation}
\text{bias} = \frac{-\lambda\beta}{1+\lambda}, \qquad \Var = \frac{\sigma^2}{(1+\lambda)^2},
\qquad
\text{MSE}(\lambda) = \frac{\sigma^2 + \lambda^2\beta^2}{(1+\lambda)^2} .
\end{equation}
Differentiating, $\text{MSE}'(0) = -2\sigma^2 < 0$: \textbf{some positive $\lambda$ always improves on OLS}, for any $\beta$ and any $\sigma^2 > 0$ (the general version is Hoerl \& Kennard's 1970 existence theorem). Setting the derivative to zero gives
\begin{equation}
\lambda^\star = \frac{\sigma^2}{\beta^2},
\qquad
\text{MSE}(\lambda^\star) = \frac{\sigma^2\beta^2}{\beta^2 + \sigma^2} \;<\; \sigma^2 = \text{MSE}(\hat\beta^{\text{OLS}}),
\end{equation}
and $\lambda^\star = \sigma^2/\beta^2$ is precisely the noise-to-prior-variance ratio that the Gaussian-prior MAP reading predicts.
\end{mathresult}

This is the bias--variance trade instantiated on the simplest possible model: pay a little squared bias, collect a larger variance reduction, and the net is negative. The decomposition itself is canonical in [DL 2]; what is worth carrying from here is the shape of the argument, because it is the same one that justifies shrinkage everywhere---smoothing in $k$-NN, shallow trees in boosting, and weight decay in a network.

\subsection{Lasso and the Constrained View}

The lasso replaces the $\ell_2$ penalty with $\ell_1$:
\begin{equation}
\hat\beta^{\text{lasso}} = \argmin_\beta \ \tfrac{1}{2}\norm{y - X\beta}_2^2 + \lambda\norm{\beta}_1 .
\end{equation}
There is no closed form (except in the orthonormal case, below), the objective is convex but not strictly convex when $d>n$ (so the coefficient vector may not be unique even though the fitted values are), and the standard solver is cyclic coordinate descent with warm starts along a $\lambda$ path---which is what \texttt{glmnet} and scikit-learn's \texttt{Lasso}/\texttt{ElasticNet} do.

The bridge to the geometry answer is \textbf{Lagrangian duality}: for every $\lambda \ge 0$ there is a $t \ge 0$ (and conversely) such that the penalized problem
\begin{equation}
\min_\beta \ \tfrac{1}{2}\norm{y-X\beta}_2^2 + \lambda\, P(\beta)
\qquad\text{is equivalent to}\qquad
\min_{\beta:\ P(\beta) \le t} \ \tfrac{1}{2}\norm{y-X\beta}_2^2 .
\end{equation}
Increasing $\lambda$ corresponds to decreasing $t$. This equivalence is the entire reason the picture in Figure~\ref{cml:fig:l1l2geometry} is legitimate: the penalized estimator really is the point where the smallest RSS contour first touches the constraint set.

\begin{figure}[H]
\centering
\begin{tikzpicture}[scale=0.9]
\begin{scope}
  \draw[->,gray!70] (-1.7,0) -- (3.5,0) node[right,black,font=\small] {$\beta_1$};
  \draw[->,gray!70] (0,-1.7) -- (0,3.5) node[above,black,font=\small] {$\beta_2$};
  \fill[blue!12,draw=deepblue,thick] (1.2,0) -- (0,1.2) -- (-1.2,0) -- (0,-1.2) -- cycle;
  \foreach \s in {0.42, 0.72, 1.0}
    \draw[deepred!75, rotate around={19.3:(1.85,2.35)}]
      (1.85,2.35) ellipse [x radius={2.253*\s}, y radius={1.428*\s}];
  \fill[deepred] (1.85,2.35) circle (1.7pt);
  \node[font=\scriptsize, black, above right] at (1.9,2.4) {$\hat\beta^{\text{OLS}}$};
  \fill[deepblue] (0,1.2) circle (2.3pt);
  \node[font=\scriptsize, black, left] at (-0.1,1.32) {$\hat\beta^{\text{lasso}}$};
  \node[font=\small] at (0.9,-2.35) {L1 ball: $\abs{\beta_1}+\abs{\beta_2} \le t$};
  \node[font=\scriptsize, deepblue, align=center, text width=4.3cm] at (0.9,-3.3) {corners on the axes\\$\Rightarrow$ $\hat\beta_1 = 0$ exactly};
\end{scope}
\begin{scope}[xshift=7.2cm]
  \draw[->,gray!70] (-1.7,0) -- (3.5,0) node[right,black,font=\small] {$\beta_1$};
  \draw[->,gray!70] (0,-1.7) -- (0,3.5) node[above,black,font=\small] {$\beta_2$};
  \fill[blue!12,draw=deepblue,thick] (0,0) circle (1.2);
  \foreach \s in {0.42, 0.72, 1.0}
    \draw[deepred!75, rotate around={19.3:(1.85,2.35)}]
      (1.85,2.35) ellipse [x radius={2.067*\s}, y radius={1.310*\s}];
  \fill[deepred] (1.85,2.35) circle (1.7pt);
  \node[font=\scriptsize, black, above right] at (1.9,2.4) {$\hat\beta^{\text{OLS}}$};
  \fill[deepblue] (0.48,1.09) circle (2.3pt);
  \node[font=\scriptsize, black, right, fill=white, inner sep=1pt] at (0.6,0.90) {$\hat\beta^{\text{ridge}}$};
  \node[font=\small] at (0.9,-2.35) {L2 ball: $\beta_1^2+\beta_2^2 \le t$};
  \node[font=\scriptsize, deepblue, align=center, text width=4.3cm] at (0.9,-3.3) {smooth everywhere\\$\Rightarrow$ both coordinates shrink, neither vanishes};
\end{scope}
\end{tikzpicture}
\caption{The constrained view of regularization. Red ellipses are contours of the residual sum of squares centered at $\hat\beta^{\text{OLS}}$; the blue region is the feasible set. The estimator is the first point of contact. The L1 ball is non-differentiable at its vertices, which lie \emph{on the coordinate axes}, and generic contours hit a vertex with positive probability---so a coefficient is set to exactly zero. The L2 ball is smooth, so the contact point almost surely has every coordinate nonzero.}
\label{cml:fig:l1l2geometry}
\end{figure}

The geometric argument in words: the constraint set for L1 is a cross-polytope whose \emph{corners lie on the coordinate axes}, and a corner is where a whole cone of contour orientations makes first contact---so hitting one is not a measure-zero coincidence but the typical outcome. Sparsity is a consequence of the \emph{shape} of the ball (non-differentiability at points where coordinates vanish), not of the penalty's \emph{strength}. This distinction is the difference between a passing and a failing answer.

\subsection{Soft-Thresholding: the Same Answer with Algebra}

The picture alone is a partial answer. The complete one adds the orthonormal-design case, where both estimators have closed forms and the mechanism becomes visible.

\begin{mathresult}{Ridge Scales, Lasso Subtracts}
Assume the columns of $X$ are orthonormal, $X^{\top}X = I$, and write $\hat\beta^{\text{OLS}} = X^{\top}y$. Then the objective separates across coordinates:
\begin{equation}
\tfrac{1}{2}\norm{y - X\beta}_2^2 + \lambda P(\beta)
= \tfrac{1}{2}\norm{\beta - \hat\beta^{\text{OLS}}}_2^2 + \lambda P(\beta) + \text{const},
\end{equation}
so each $\beta_j$ solves a one-dimensional problem. Minimizing $\tfrac{1}{2}(\beta_j - b_j)^2 + \tfrac{\lambda}{2}\beta_j^2$ and $\tfrac{1}{2}(\beta_j - b_j)^2 + \lambda\abs{\beta_j}$ respectively gives
\begin{align}
\textbf{Ridge:}\quad & \hat\beta_j = \frac{\hat\beta_j^{\text{OLS}}}{1+\lambda} && \text{(uniform multiplicative shrinkage)} \\
\textbf{Lasso:}\quad & \hat\beta_j = \operatorname{sign}\!\big(\hat\beta_j^{\text{OLS}}\big)\Big(\big|\hat\beta_j^{\text{OLS}}\big| - \lambda\Big)_{+} && \text{(\textbf{soft-thresholding})}
\end{align}
The lasso step comes from the subgradient condition: for $\beta_j \ne 0$, $\beta_j - b_j + \lambda\operatorname{sign}(\beta_j) = 0$; for $\beta_j = 0$ the subdifferential of $\abs{\cdot}$ is the whole interval $[-\lambda,\lambda]$, so $\beta_j = 0$ is optimal whenever $\abs{b_j} \le \lambda$. \textbf{That interval is the sparsity mechanism.} Ridge's derivative $\lambda\beta_j$ vanishes as $\beta_j \to 0$, so it can never push a coefficient the last inch to zero; the lasso's subgradient has constant magnitude $\lambda$ all the way in.
\end{mathresult}

Worked numerically, because this is what convinces an interviewer you own it. Take $\hat\beta^{\text{OLS}} = (2.0,\ 0.4,\ -0.9)$ and $\lambda = 0.5$:
\begin{itemize}
    \item \textbf{Lasso:} $(2.0 - 0.5,\ 0,\ -(0.9-0.5)) = (1.5,\ \mathbf{0},\ -0.4)$. The small coefficient is annihilated; the survivors are shifted toward zero by a \emph{constant}, which is also why the lasso is biased for large coefficients (the motivation for the relaxed lasso, adaptive lasso, and SCAD/MCP).
    \item \textbf{Ridge:} $(2.0,\ 0.4,\ -0.9)/1.5 = (1.333,\ 0.267,\ -0.6)$. Every coefficient shrinks by the same \emph{factor}; none reaches zero.
\end{itemize}

The orthonormal case is not a toy: the soft-threshold operator \emph{is} the coordinate-descent update in the general case. With standardized columns, cycling
\begin{equation}
\beta_j \leftarrow S_\lambda\!\left(x_j^{\top}\big(y - X_{-j}\beta_{-j}\big)\right), \qquad S_\lambda(z) = \operatorname{sign}(z)(\abs{z}-\lambda)_+ ,
\end{equation}
to convergence is how production lasso solvers work---each coordinate faces exactly the one-dimensional problem above, with the other coordinates' fit removed. Being able to say this closes the loop between the picture, the algebra, and the implementation.

\subsection{The Regularization Path}

Fitting at a single $\lambda$ is rarely how these models are used; the object of interest is the \textbf{path}, the whole family $\{\hat\beta(\lambda) : \lambda \ge 0\}$ traced from the null model down to OLS. Two facts make the path cheap and informative.

\textbf{The lasso path is piecewise linear in $\lambda$.} Between the ``knots'' where a variable enters or leaves the active set, the solution moves linearly, so the entire path can be computed exactly in roughly the cost of one least-squares fit (LARS; Efron et al., 2004). Modern solvers usually prefer coordinate descent on a log-spaced grid with warm starts, which is nearly as cheap and generalizes to logistic and other losses; either way, ``fit the path, then pick $\lambda$'' costs about what one fit costs. The ridge path, by contrast, is smooth in $\lambda$---no variable ever enters or leaves, because nothing is ever exactly zero.

\textbf{Read the path, do not just read the winner.} The order in which variables enter is a rough signal of strength; a coefficient that changes sign along the path is a collinearity warning; and a validation curve that is flat over an order of magnitude of $\lambda$ says the regularization strength is not the binding constraint, so tuning it further is wasted effort. This is also where the one-standard-error rule is applied---on the curve, not on a single number.

\subsection{The Bayesian Reading, and Elastic Net}

Both penalties are MAP estimates under a prior on $\beta$: L2 corresponds to a Gaussian prior $\beta_j \sim \mathcal{N}(0, \tau^2)$ with $\lambda = \sigma^2/\tau^2$, and L1 to a Laplace prior $\beta_j \sim \mathrm{Laplace}(0,b)$ with $\lambda = \sigma^2/b$. The Laplace density has a spike at the origin, which is the probabilistic restatement of the corner. The full MAP machinery---and the important caveat that the lasso's \emph{point estimate} is the Laplace MAP but the Laplace \emph{posterior mean} is not sparse---belongs to the probabilistic-foundations chapter (Chapter~\ref{cml:chap:probabilistic_foundations}); here it is a pointer and a one-line answer to ``is there a Bayesian interpretation?''

The \textbf{elastic net} combines them,
\begin{equation}
\min_\beta \ \tfrac{1}{2}\norm{y-X\beta}_2^2 + \lambda\Big(\alpha\norm{\beta}_1 + \tfrac{1-\alpha}{2}\norm{\beta}_2^2\Big),
\end{equation}
and exists to fix two concrete lasso pathologies. First, with a group of highly correlated predictors the lasso picks essentially one of them \emph{arbitrarily}---which member survives is decided by noise, so the selected set is unstable across bootstrap resamples and useless as a scientific claim. The ridge component induces a \textbf{grouping effect}: correlated predictors receive similar coefficients and enter or leave together. Second, when $d > n$ the lasso saturates at $n$ selected variables; the elastic net has no such ceiling. The price is a second hyperparameter, and the standard recipe is a small grid over $\alpha$ crossed with a warm-started path over $\lambda$.

\subsection{Practicalities That Get Candidates Caught}

\begin{warningbox}
\textbf{Standardize first, and do not penalize the intercept.} Penalties are not scale invariant: rescale $x_j$ by $c$ and $\hat\beta_j$ rescales by $1/c$, so $\lambda\abs{\beta_j}$ changes---an unstandardized penalty silently ranks features by their measurement units, and a feature in dollars is penalized a thousand times more heavily than the same feature in thousands of dollars. Penalizing the intercept is worse than useless: it makes predictions depend on where you put the origin of $y$, so shifting all targets by a constant changes the fit. The standard implementation centers $y$ and $X$, scales $X$ to unit variance, fits without an intercept, and back-transforms. Standardization must be fit on training folds only---the scaler is part of the pipeline, and leaking it is a classic cross-validation bug (Chapter~\ref{cml:chap:applied_craft}).
\end{warningbox}

Three more items worth having ready:
\begin{itemize}
    \item \textbf{Choosing $\lambda$.} Cross-validation over a log-spaced grid, using the full path (warm starts make the whole path cost roughly one fit). The \textbf{one-standard-error rule}---take the largest $\lambda$ whose CV error is within one standard error of the minimum---buys a simpler, more stable model at negligible cost and is a good thing to volunteer.
    \item \textbf{Lasso for selection is not lasso for inference.} Post-selection coefficients and their $p$-values are invalid if computed as though the feature set were chosen a priori; the honest routes are stability selection, the debiased/desparsified lasso, or split-sample selection and inference.
    \item \textbf{Weight decay in deep learning is not always L2.} With plain SGD they coincide, but with adaptive optimizers the L2 term is rescaled by the per-parameter second-moment estimate and the equivalence breaks---which is what AdamW's decoupling restores. That interaction is canonical in [DL 15]; the classical statement here is the base case it decouples from.
\end{itemize}

\begin{table}[H]
\centering\small
\setlength{\tabcolsep}{4pt}
\caption{The three penalties, side by side.}
\label{cml:tab:penalties}
\begin{tabularx}{\textwidth}{@{}>{\raggedright\arraybackslash}p{2.45cm} X X X@{}}
\toprule
& \textbf{Ridge (L2)} & \textbf{Lasso (L1)} & \textbf{Elastic net} \\
\midrule
Penalty & $\tfrac{\lambda}{2}\norm{\beta}_2^2$ & $\lambda\norm{\beta}_1$ & both, mixed by $\alpha$ \\
Closed form & yes, $(X^{\top}X{+}\lambda I)^{-1}X^{\top}y$ & no & no \\
Orthonormal case & scale by $1/(1{+}\lambda)$ & soft-threshold at $\lambda$ & threshold, then scale \\
Sparse? & never & yes & yes \\
Unique solution & always & not when $d>n$ & always ($\alpha<1$) \\
Correlated group & splits weight evenly & picks one arbitrarily & grouping effect \\
Cap on selected & --- & at most $n$ when $d>n$ & none \\
Prior (MAP) & Gaussian & Laplace & Gaussian $\times$ Laplace \\
Reach for it when & collinearity; every feature plausibly matters; $d>n$ & selection wanted; a small serving model & correlated groups; $d \gg n$ \\
\bottomrule
\end{tabularx}
\end{table}

%==============================================================================
\section{Logistic Regression and the GLM View}
\label{cml:sec:logistic_glm}
%==============================================================================

\subsection{From the Bernoulli Likelihood to Cross-Entropy}

Do not start from ``we squash the output with a sigmoid.'' Start from the modeling assumption, because that is what the derivation question is checking.

\begin{mathresult}{Logistic Regression, Derived}
Assume $y_i \in \{0,1\}$ with $y_i \given x_i \sim \mathrm{Bernoulli}(p_i)$, and model the \textbf{log-odds} as linear:
\begin{equation}
\log\frac{p_i}{1-p_i} = x_i^{\top}\beta
\qquad\Longleftrightarrow\qquad
p_i = \sigma(x_i^{\top}\beta) = \frac{1}{1+e^{-x_i^{\top}\beta}} .
\end{equation}
The sigmoid is not an arbitrary squashing choice---it is the \emph{inverse of the assumed link}. The likelihood and its negative log are
\begin{equation}
L(\beta) = \prod_{i=1}^{n} p_i^{y_i}(1-p_i)^{1-y_i},
\qquad
\loss(\beta) = -\sum_{i=1}^{n}\Big[y_i\log p_i + (1-y_i)\log(1-p_i)\Big],
\end{equation}
which is exactly binary \textbf{cross-entropy}: log loss is not a separate design decision, it \emph{is} the Bernoulli negative log-likelihood. Differentiating, and using $\sigma'(z) = \sigma(z)(1-\sigma(z))$:
\begin{equation}
\frac{\partial \loss}{\partial \beta} = \sum_{i=1}^{n}\big(\sigma(x_i^{\top}\beta) - y_i\big)\,x_i
\qquad\Longleftrightarrow\qquad
\boxed{\ \nabla_\beta \loss = X^{\top}\big(\sigma(X\beta) - y\big)\ }
\end{equation}
The Hessian is
\begin{equation}
\nabla^2_\beta \loss = X^{\top}W X, \qquad W = \diag\big(p_i(1-p_i)\big) \succeq 0 ,
\end{equation}
so $\loss$ is convex---every local minimum is global---and strictly convex when $X$ has full column rank and no $p_i$ is at $0$ or $1$.
\end{mathresult}

The two-line derivation of the gradient, which is worth being able to produce term by term:
$\partial \loss_i/\partial p_i = -(y_i/p_i - (1-y_i)/(1-p_i)) = (p_i - y_i)/(p_i(1-p_i))$, and $\partial p_i/\partial z_i = p_i(1-p_i)$ where $z_i = x_i^{\top}\beta$. The $p_i(1-p_i)$ factors cancel exactly, leaving $\partial \loss_i/\partial z_i = p_i - y_i$. That cancellation is the whole reason the gradient is clean, and it is the first thing Section~\ref{cml:sec:glm_family} generalizes.

\begin{keyinsight}[Same Gradient Shape as Linear Regression]
Least squares gives $\nabla = X^{\top}(X\beta - y)$; logistic regression gives $\nabla = X^{\top}(\sigma(X\beta) - y)$. Both are $X^{\top}(\hat{y} - y)$: \textbf{design matrix transpose times the prediction error}. This is not a coincidence and not a mnemonic---it is a theorem about exponential-family models fitted through their canonical link (Section~\ref{cml:sec:glm_family}). Softmax regression, Poisson regression, and the residual definition inside gradient boosting (Chapter~\ref{cml:chap:trees_ensembles}) are all the same statement.
\end{keyinsight}

\subsection{No Closed Form, and How It Is Actually Fit}

Setting $X^{\top}(\sigma(X\beta) - y) = 0$ gives $d$ equations that are transcendental in $\beta$; there is no algebraic solution, and ``logistic regression has a closed form'' is a hard red flag. What exists instead is a very well-behaved convex optimization problem:

\begin{itemize}
    \item \textbf{Newton / IRLS} (the classical method). The Newton step rearranges into a \emph{weighted least squares} fit:
    \begin{equation*}
    \beta^{(t+1)} = \beta^{(t)} + (X^{\top}WX)^{-1}X^{\top}(y - p) = (X^{\top}WX)^{-1}X^{\top}W z,
    \end{equation*}
    with working response $z = X\beta^{(t)} + W^{-1}(y-p)$---hence \emph{iteratively reweighted least squares}. It converges quadratically in a handful of iterations but costs $O(nd^2 + d^3)$ per step, so it is the right choice only for moderate $d$. This is what R's \texttt{glm} does.
    \item \textbf{Quasi-Newton} (L-BFGS): the default in scikit-learn. It stores a low-rank Hessian approximation, costing $O(nd)$ per iteration with near-Newton convergence.
    \item \textbf{SGD / coordinate descent}: for very large $n$ or $d$, or for L1 penalties. FTRL-Proximal is the industrial online variant (Section~\ref{cml:sec:supervised_interview}).
\end{itemize}

\subsection{Why Log Loss and Not Squared Error}

A favorite probe, and a place where a shallow answer (``because it is a classification problem'') is easy to spot. There are three distinct reasons, and the strong answer separates them.

\textbf{1. Maximum likelihood.} Squared error is the Gaussian NLL; the label is Bernoulli. Log loss is the likelihood the model actually implies, and it makes the estimator consistent and asymptotically efficient under that model.

\textbf{2. Convexity.} Squared error composed with a sigmoid is \emph{not convex} in $\beta$, so gradient descent can stall at a bad stationary point. The counterexample is small enough to give on a whiteboard: with one feature $x=1$ and label $y=0$, the objective is $f(\beta) = \tfrac12\sigma(\beta)^2$, and
\begin{equation}
f''(\beta) = \sigma^2(1-\sigma)(2 - 3\sigma) < 0 \quad\text{whenever}\quad \sigma(\beta) > \tfrac{2}{3},
\end{equation}
i.e., for $\beta > \log 2 \approx 0.69$. Log loss, by contrast, has Hessian $X^{\top}WX \succeq 0$ everywhere.

\textbf{3. Gradient saturation.} Differentiating the two losses with respect to the logit $z$:
\begin{equation}
\frac{\partial}{\partial z}\,\tfrac12(\sigma(z)-y)^2 = (\sigma(z)-y)\,\sigma(z)(1-\sigma(z)),
\qquad
\frac{\partial}{\partial z}\,\text{logloss} = \sigma(z)-y .
\end{equation}
The squared-error gradient carries an extra $\sigma'(z)$ factor that vanishes when the model is confident. Concretely, at $z = -10$ with true label $y=1$---confidently, catastrophically wrong---the log-loss gradient is $\sigma(-10)-1 \approx -0.99995$ while the squared-error gradient is about $-4.5\times10^{-5}$: a factor of roughly $22{,}000$ smaller, so the most informative example in the batch contributes almost nothing. Log loss is unbounded in exactly the region where the model deserves a large correction.

\begin{warningbox}
The one nuance that separates a good answer from a great one: \textbf{squared error on probabilities is the Brier score, which is a proper scoring rule}, so ``MSE gives uncalibrated probabilities'' is wrong. Both losses are minimized in expectation at the true $\Pr(y=1\given x)$. The case against squared error here is about the \emph{likelihood justification} and the \emph{optimization geometry}, not about propriety. Candidates who assert MSE is improper are confidently wrong about a fact the interviewer may well know; candidates who name the Brier score and then give the convexity and saturation arguments are visibly ahead.
\end{warningbox}

\subsection{Separation: When the Weights Blow Up}

If some hyperplane perfectly separates the classes, then for any separating $\beta$ the scaled vector $c\beta$ drives every $p_i$ toward the correct $0$ or $1$ as $c \to \infty$, so the likelihood increases monotonically and \textbf{the MLE does not exist}---the optimizer runs to infinity until it hits an iteration cap or numerical limits. \emph{Quasi}-complete separation (perfect separation within a subset, typically caused by a rare category that appears with only one label) produces the same divergence in a subset of coordinates.

Symptoms: coefficients of magnitude $10^3$--$10^6$, standard errors even larger, training accuracy exactly $1.0$, fitted probabilities pinned at $0$ or $1$, and convergence warnings. Fixes, in order of preference:
\begin{itemize}
    \item \textbf{Any} penalty. L2 makes the objective strongly convex, so the optimum is finite and unique; this is why scikit-learn regularizes logistic regression \emph{by default} (and why a candidate who has never seen the divergence is often just someone who never turned the penalty off).
    \item \textbf{Firth's penalized likelihood} (equivalently a Jeffreys prior), the standard fix in the statistics literature, which also removes small-sample bias.
    \item \textbf{Investigate the feature}, because separation is frequently a \textbf{leakage alarm}: a feature that perfectly predicts the label is more often a post-outcome variable than a discovery. Check the timestamps before celebrating.
    \item Collapse or smooth the offending rare category; note that dropping the feature is a last resort, not a diagnosis.
\end{itemize}

\subsection{The GLM Frame}
\label{cml:sec:glm_family}

Generalized linear models are the structure that makes linear, logistic, Poisson, and softmax regression one object rather than four tricks. A one-parameter exponential-family density in canonical form is
\begin{equation}
f(y \given \theta, \phi) = \exp\left\{\frac{y\theta - b(\theta)}{a(\phi)} + c(y,\phi)\right\},
\end{equation}
with the two facts that make the whole machinery run: $\E[y] = \mu = b'(\theta)$ and $\Var(y) = a(\phi)\,b''(\theta)$. A GLM adds a \textbf{link} $g$ with $g(\mu) = x^{\top}\beta$. The link is \textbf{canonical} when it maps the mean straight onto the natural parameter, $\theta = x^{\top}\beta$.

\begin{mathresult}{Why the Canonical Link Gives the $(\hat{y}-y)$ Gradient}
With the canonical link, $\theta_i = x_i^{\top}\beta$, so the log-likelihood contribution is $\ell_i = \{y_i x_i^{\top}\beta - b(x_i^{\top}\beta)\}/a(\phi)$ and
\begin{equation}
\frac{\partial \ell_i}{\partial \beta} = \frac{y_i - b'(x_i^{\top}\beta)}{a(\phi)}\, x_i = \frac{y_i - \mu_i}{a(\phi)}\,x_i
\qquad\Longrightarrow\qquad
\nabla_\beta \loss = \frac{1}{a(\phi)} X^{\top}(\mu - y).
\end{equation}
The identity $b'(\theta) = \mu$ is what makes the awkward chain-rule factor cancel. So the ``prediction minus target'' gradient is a \emph{property of the canonical link}, not a lucky feature of the sigmoid---and it also explains why the Fisher information equals the observed Hessian there, which is why IRLS and Fisher scoring coincide for canonical links.
\end{mathresult}

\begin{table}[H]
\centering\small
\caption{The GLM family. Every row is fitted by the same IRLS loop; only $b(\theta)$ changes.}
\label{cml:tab:glm}
\begin{tabularx}{\textwidth}{@{}l l l l X@{}}
\toprule
\textbf{Response} & \textbf{Family} & \textbf{Canonical link} & \textbf{$b(\theta)$} & \textbf{Mean/variance and typical use} \\
\midrule
Continuous & Gaussian & identity & $\theta^2/2$ & $\Var = \sigma^2$, constant. Ordinary regression. \\
Binary & Bernoulli & logit & $\log(1+e^\theta)$ & $\Var=\mu(1-\mu)$. Classification, CTR. \\
Counts & Poisson & log & $e^\theta$ & $\Var=\mu$. Events per exposure; use an offset $\log(\text{exposure})$. \\
Positive skewed & Gamma & inverse & $-\log(-\theta)$ & $\Var=\phi\mu^2$, constant CV. Durations, claim sizes. \\
$K$ classes & Multinomial & softmax & --- & Softmax regression. \\
\bottomrule
\end{tabularx}
\end{table}

Two GLM facts that come up as follow-ups. \textbf{Overdispersion}: the Poisson family forces $\Var = \mu$, and real count data routinely has $\Var \gg \mu$; the symptom is a residual deviance far above the residual degrees of freedom, and the fixes are quasi-Poisson (inflate the standard errors by an estimated dispersion) or negative binomial (add a genuine dispersion parameter). \textbf{Log link versus log target}: modeling $\E[\log y]$ with OLS is not the same as modeling $\log \E[y]$ with a Poisson/Gamma GLM---the first estimates a median-like quantity on the original scale and cannot handle $y=0$, and back-transforming needs a smearing correction. Choosing the GLM is usually the right answer for positive, skewed, or count targets.

\subsection{Multiclass, and Reading Coefficients}

\textbf{Softmax regression} extends the same derivation: $p_k(x) = e^{x^{\top}\beta_k}/\sum_{k'} e^{x^{\top}\beta_{k'}}$, with NLL gradient $\nabla_{\beta_k} = X^{\top}(p_k - \mathds{1}[y=k])$---the same shape again. The parameterization is \textbf{overcomplete}: adding any fixed vector $c$ to every $\beta_k$ leaves all probabilities unchanged, so only $K-1$ blocks are identifiable. Statistics packages resolve this by fixing a reference class ($\beta_K = 0$); machine-learning libraries typically fit all $K$ blocks and let the L2 penalty pick the unique minimum-norm representative---which is a fine engineering answer but means the raw coefficients are not comparable to a statistics package's without care. \textbf{One-vs-rest} trains $K$ independent binary models; it is simpler and parallel, but the $K$ scores are not jointly normalized, so its probabilities need rescaling and it can behave poorly under class imbalance. Softmax is the default when the classes are genuinely mutually exclusive.

\textbf{Interpretation.} $\beta_j$ is the additive change in log-odds per unit increase in $x_j$, all else held fixed; $e^{\beta_j}$ is the corresponding \emph{odds ratio}. Three traps: (i) comparing raw coefficient magnitudes across features with different units is meaningless---standardize first, and even then the comparison is about association strength, not importance in any causal sense; (ii) odds ratios are not risk ratios, and the gap is large when the base rate is high; (iii) ``all else held fixed'' is a statement about the model, not the world, and collinearity makes individual coefficients unstable (Section~\ref{cml:sec:linear_regression}) even when the fit is excellent.

%==============================================================================
\section{Support Vector Machines and Kernels}
\label{cml:sec:svm_kernels}
%==============================================================================

\subsection{The Margin, Derived}

Let $f(x) = w^{\top}x + b$ and $y_i \in \{-1,+1\}$. The signed distance from $x_i$ to the hyperplane $f(x)=0$ is $(w^{\top}x_i+b)/\norm{w}$, so the \textbf{geometric margin} of the training set is $\min_i y_i(w^{\top}x_i+b)/\norm{w}$.

The scaling $(w,b) \mapsto (cw, cb)$ leaves the hyperplane and every geometric margin unchanged, so the problem as stated is underdetermined. Fix it by \textbf{canonical scaling}: require the closest points to satisfy $\min_i y_i(w^{\top}x_i+b) = 1$. The geometric margin is then exactly $1/\norm{w}$, the separating band has width $2/\norm{w}$, and maximizing the margin becomes

\begin{mathresult}{Hard-Margin Primal}
\begin{equation}
\min_{w,b}\ \tfrac{1}{2}\norm{w}_2^2 \qquad \text{subject to}\qquad y_i(w^{\top}x_i + b) \ge 1 \quad \forall i .
\end{equation}
A convex quadratic program: quadratic objective, linear constraints, hence a unique global optimum. Maximizing $2/\norm{w}$ and minimizing $\tfrac12\norm{w}^2$ are the same problem; the square and the $\tfrac12$ are conveniences for differentiation.
\end{mathresult}

Real data is rarely separable, and even when it is, a hard margin is a variance disaster (one mislabeled point can swing the boundary arbitrarily). Introduce slack $\xi_i \ge 0$:
\begin{equation}
\min_{w,b,\xi}\ \tfrac{1}{2}\norm{w}_2^2 + C\sum_{i=1}^n \xi_i
\qquad\text{s.t.}\qquad y_i(w^{\top}x_i+b) \ge 1 - \xi_i, \quad \xi_i \ge 0 .
\end{equation}

\begin{keyinsight}[Constrained Form and Loss Form Are the Same Object]
At the optimum each $\xi_i$ takes its smallest feasible value, $\xi_i = \max(0,\,1 - y_i f(x_i))$. Substituting eliminates the constraints entirely:
\begin{equation}
\min_{w,b}\ \tfrac{1}{2}\norm{w}_2^2 + C\sum_{i=1}^n \max\big(0,\ 1 - y_i f(x_i)\big) ,
\end{equation}
which after dividing by $C$ is ordinary \textbf{regularized empirical risk minimization with hinge loss} and $\lambda = 1/(2C)$. Moving fluidly between the two forms is the thing the SVM question is really testing: the geometric story (maximize the margin) and the statistical story (hinge loss plus L2 penalty) are one derivation apart, and the second is what tells you $C$ is an inverse regularization strength.
\end{keyinsight}

\begin{figure}[H]
\centering
\begin{tikzpicture}
\begin{axis}[width=0.80\textwidth, height=5.8cm,
  xlabel={margin $m = y\,f(x)$}, ylabel={loss},
  xmin=-2.5, xmax=2.5, ymin=-0.15, ymax=3.7,
  legend pos=north east, legend style={font=\scriptsize, draw=gray!50},
  grid=major, grid style={gray!20},
  tick label style={font=\scriptsize}, label style={font=\small}]
\addplot[deepblue, very thick, domain=-2.5:2.5, samples=200] {max(0,1-x)};
\addlegendentry{hinge $\max(0,1-m)$}
\addplot[deepred, very thick, domain=-2.5:2.5, samples=200] {ln(1+exp(-x))/ln(2)};
\addlegendentry{logistic $\log_2(1+e^{-m})$}
\addplot[deepgreen, thick, dashed, domain=-2.5:2.5, samples=200] {min(3.7,(1-x)^2)};
\addlegendentry{squared $(1-m)^2$}
\addplot[gray, thick, dotted] coordinates {(-2.5,1) (0,1) (0,0) (2.5,0)};
\addlegendentry{0--1 loss}
\end{axis}
\end{tikzpicture}
\caption{Surrogate losses as functions of the margin $m = y f(x)$. All three upper-bound the 0--1 loss and are convex. Hinge is \emph{exactly} zero for $m \ge 1$---the source of the SVM's sparse solution, since correctly classified interior points contribute nothing. Logistic is strictly positive everywhere, so every point keeps a small pull, which is why logistic regression has no support vectors but does produce probabilities. Squared loss \emph{penalizes being too correct} ($m > 1$), which is why it is a poor classification surrogate.}
\label{cml:fig:margin_losses}
\end{figure}

\subsection{The Dual, KKT, and What a Support Vector Is}

The dual is where the kernel trick comes from, so the derivation is worth doing rather than asserting. Form the Lagrangian with multipliers $\alpha_i \ge 0$ for the margin constraints and $\mu_i \ge 0$ for $\xi_i \ge 0$:
\begin{equation}
\mathcal{L} = \tfrac12\norm{w}^2 + C\sum_i \xi_i - \sum_i \alpha_i\big[y_i(w^{\top}x_i+b) - 1 + \xi_i\big] - \sum_i \mu_i \xi_i .
\end{equation}
Stationarity in the primal variables gives three conditions:
\begin{equation}
\frac{\partial \mathcal{L}}{\partial w} = 0 \Rightarrow w = \sum_i \alpha_i y_i x_i,
\qquad
\frac{\partial \mathcal{L}}{\partial b} = 0 \Rightarrow \sum_i \alpha_i y_i = 0,
\qquad
\frac{\partial \mathcal{L}}{\partial \xi_i} = 0 \Rightarrow \alpha_i + \mu_i = C .
\end{equation}
The first is the \textbf{representer} statement---the optimal weight vector is a linear combination of training points---and the third, with $\mu_i \ge 0$, is what turns the multiplier constraint into the \textbf{box} $0 \le \alpha_i \le C$. Substituting back eliminates $w$, $b$, and $\xi$ entirely:

\begin{mathresult}{The Dual QP}
\begin{equation}
\max_{\alpha}\ \sum_{i=1}^n \alpha_i - \tfrac12 \sum_{i=1}^n\sum_{j=1}^n \alpha_i\alpha_j\, y_iy_j\, \langle x_i, x_j\rangle
\qquad\text{s.t.}\qquad 0 \le \alpha_i \le C,\quad \sum_i \alpha_i y_i = 0 .
\end{equation}
The decision function is likewise expressible in inner products only:
\begin{equation}
f(x) = \sum_{i=1}^{n} \alpha_i y_i \langle x_i, x\rangle + b .
\end{equation}
\textbf{The training data enters only through pairwise inner products.} That single observation is the kernel trick's entire precondition.
\end{mathresult}

Complementary slackness, $\alpha_i[y_i f(x_i) - 1 + \xi_i] = 0$ and $\mu_i\xi_i = 0$, partitions the training set into exactly three cases---this is the crisp answer to ``what is a support vector?''

\begin{table}[H]
\centering\small
\caption{Reading the training set off the KKT conditions.}
\label{cml:tab:svm_kkt}
\begin{tabularx}{\textwidth}{@{}l l l X@{}}
\toprule
\textbf{Multiplier} & \textbf{Slack} & \textbf{Position} & \textbf{Role} \\
\midrule
$\alpha_i = 0$ & $\xi_i = 0$ & $y_i f(x_i) > 1$, strictly outside the margin & Not a support vector; deleting it does not change the solution. \\
$0 < \alpha_i < C$ & $\xi_i = 0$ & $y_i f(x_i) = 1$, exactly on the margin & \emph{Free} (margin) support vector. Used to recover $b = y_i - \sum_j \alpha_j y_j K(x_j,x_i)$. \\
$\alpha_i = C$ & $\xi_i \ge 0$ & $y_i f(x_i) \le 1$, inside the margin or misclassified & \emph{Bounded} support vector. \\
\bottomrule
\end{tabularx}
\end{table}

Two consequences worth stating unprompted. The solution is \textbf{sparse in examples}: only support vectors carry nonzero $\alpha_i$, so inference costs $O(n_{\text{SV}} \cdot d)$ rather than $O(n \cdot d)$ (contrast the lasso, which is sparse in \emph{features}). And the SVM's boundary depends only on points near it---the interior of each class could be resampled freely without changing the model, which is simultaneously its robustness story and the reason it is exquisitely sensitive to mislabeled points near the boundary.

\textbf{Primal or dual?} Solve the \emph{dual} when you want a kernel (it is the only form that admits one) or when $n \ll d$, since the dual has $n$ variables against the primal's $d$. Solve the \emph{primal} when $n \gg d$ and a linear kernel suffices---which is exactly the regime LIBLINEAR and SGD-on-hinge target, and why ``linear SVM'' and ``kernel SVM'' behave like different tools with different scaling laws despite sharing a derivation.

\subsection{Kernels}

The data appears only as $\langle x_i, x_j\rangle$. Replace that inner product with $K(x_i,x_j) = \langle \phi(x_i), \phi(x_j)\rangle$ for some feature map $\phi$ into a possibly infinite-dimensional space, and you have fitted a linear model in $\phi$-space without ever visiting it.

\begin{warningbox}
The kernel trick does \textbf{not} ``compute the high-dimensional features efficiently.'' It never computes them \emph{at all}. $\phi(x)$ is never formed, never stored, and for the RBF kernel could not be---it is infinite-dimensional. What is computed is a scalar per pair. Getting this wrong is one of the most reliable red flags on the topic, because it reveals the candidate learned the slogan rather than the dual.
\end{warningbox}

\textbf{When is a $K$ a valid kernel?} \textbf{Mercer's condition}: $K$ must be symmetric and positive semi-definite, meaning the Gram matrix $[K(x_i,x_j)]_{ij}$ is PSD for every finite set of points. PSD-ness is what guarantees a feature map exists (and that the dual QP stays convex). Kernels are closed under addition, positive scaling, products, and composition with a positive power series, which is how new kernels are built. The $\tanh$ ``sigmoid kernel'' is the standard counterexample: it is not PSD for most parameter settings, and solvers can fail on it.

\textbf{The standard menu.} Linear $K = x^{\top}x'$; polynomial $K = (\gamma x^{\top}x' + r)^p$, whose feature map is all monomials up to degree $p$ (finite-dimensional---you could form it, you just would not want to); and the workhorse \textbf{RBF} (Gaussian)
\begin{equation}
K(x,x') = \exp\big(-\gamma\norm{x-x'}^2\big), \qquad \gamma = \frac{1}{2\sigma^2}.
\end{equation}
Its feature map is genuinely infinite-dimensional, and the two-line proof is a good thing to have ready:
\begin{equation}
e^{-\gamma\norm{x-x'}^2} = e^{-\gamma\norm{x}^2}e^{-\gamma\norm{x'}^2}\,e^{2\gamma x^{\top}x'},
\qquad
e^{2\gamma x^{\top}x'} = \sum_{k=0}^{\infty}\frac{(2\gamma)^k}{k!}\,(x^{\top}x')^k ,
\end{equation}
an infinite positively weighted sum of polynomial kernels---hence an infinite-dimensional feature space containing monomials of every degree.

\begin{keyinsight}[$C$ and $\gamma$ Are Both Bias--Variance Dials, on Different Axes]
$C$ is \textbf{inverse regularization}: large $C$ makes slack expensive, pushing toward a hard margin that contorts to fit every point (low bias, high variance); small $C$ tolerates violations, giving a wide smooth margin (high bias, low variance). $\gamma$ is the \textbf{inverse squared kernel width}: large $\gamma$ makes each training point's influence a narrow spike, so the boundary becomes wiggly and, in the limit, memorizes the training set; small $\gamma$ makes the kernel nearly constant and the model nearly linear. They interact---the standard practice is a joint log-spaced grid, typically $C \in \{10^{-2},\dots,10^{3}\}$ and $\gamma \in \{10^{-4},\dots,10^{1}\}$, with scikit-learn's default $\gamma = 1/(d\cdot \Var(X))$ as the anchor. Both must be tuned on \emph{scaled} features: RBF distances are dominated by whichever feature has the largest units, so an unscaled SVM is effectively a one-feature model. The decomposition these dials trade against is [DL 2].
\end{keyinsight}

\subsection{Costs, Probabilities, and the Honest 2026 Position}

\textbf{Scaling.} The kernel matrix is $n \times n$: $O(n^2)$ memory, and SMO-style solvers run between $O(n^2)$ and $O(n^3)$ time. At $n = 10^5$, storing the Gram matrix in float64 is $10^{10}\times 8$ bytes $= 80$\,GB---the reason kernel SVMs simply stopped being used on modern datasets. The escapes are all approximations: linear SVMs solved in the primal (LIBLINEAR, or SGD on the hinge loss, both $O(nd)$ per pass), Nystr\"om or random-Fourier-feature approximations that build an explicit low-dimensional $\hat\phi$ and then fit a linear model, or a different model family.

\textbf{No probabilities.} The SVM optimizes a margin, not a likelihood; $f(x)$ is a signed distance with no probabilistic meaning, and thresholding it at $0$ is a decision, not a probability. The standard patch is \textbf{Platt scaling}---fit $\Pr(y=1\given f) = \sigma(Af+B)$ on held-out data (scikit-learn's \texttt{probability=True} does this with internal cross-validation, at roughly $5\times$ the training cost and with occasional disagreement between \texttt{predict} and \texttt{predict\_proba}). Calibration methods and their diagnostics live in Chapter~\ref{cml:chap:probabilistic_foundations} and [DL 23].

\begin{keyinsight}[Why SVMs Are Still Asked in 2026]
Say this out loud in the interview rather than pretending SVMs are a live default. \textbf{They are rarely the production choice}: gradient-boosted trees own tabular data (Chapter~\ref{cml:chap:trees_ensembles}) and neural networks own everything with structure, while the $O(n^2)$ kernel matrix rules out modern data sizes and the absence of native probabilities is a real operational cost. They persist in interviews because the primal$\to$Lagrangian$\to$dual$\to$kernel derivation is the cleanest available test of optimization literacy---duality, KKT conditions, and the difference between a constraint and a penalty---in a setting small enough for a whiteboard. And the \emph{kernel idea} is very much alive: Gaussian processes are kernel methods with a posterior attached, attention is readable as kernel regression over the sequence [NLP 5], and neural tangent kernel theory analyses wide networks as linear models in a fixed kernel. Where a linear SVM still ships: very high-dimensional sparse text with $d \gg n$, and small-$n$ scientific settings where a margin and a good kernel beat anything that needs data.
\end{keyinsight}

For completeness at name-drop depth: \textbf{SVR} swaps hinge for the $\varepsilon$-insensitive loss (zero penalty inside a tube of width $\varepsilon$, hence sparsity in regression too); \textbf{one-class SVM} learns a boundary around the data for novelty detection; and \textbf{multiclass} SVMs are built by one-vs-one (LIBSVM's default, $K(K-1)/2$ binary problems, each small) or one-vs-rest, since the margin formulation is intrinsically binary.

%==============================================================================
\section{Naive Bayes and the Generative--Discriminative Axis}
\label{cml:sec:naive_bayes}
%==============================================================================

\subsection{The Decision Rule}

\begin{mathresult}{Naive Bayes in Three Steps}
Start from Bayes' rule and drop the evidence term, which is constant across classes:
\begin{equation}
\hat c(x) = \argmax_{c}\ \Pr(c \given x) = \argmax_c \frac{\Pr(c)\Pr(x\given c)}{\Pr(x)} = \argmax_c\ \Pr(c)\,\Pr(x_1,\dots,x_d \given c).
\end{equation}
The joint $\Pr(x_1,\dots,x_d\given c)$ needs $O(K\cdot V^d)$ parameters and is unlearnable. Impose the \textbf{naive} assumption---features conditionally independent given the class, $x_i \indep x_j \given c$---which factorizes it:
\begin{equation}
\hat c(x) = \argmax_c\ \Pr(c)\prod_{j=1}^{d}\Pr(x_j \given c)
= \argmax_c\ \left[\log \Pr(c) + \sum_{j=1}^{d}\log \Pr(x_j\given c)\right] .
\end{equation}
Parameter count collapses to $O(K\cdot V\cdot d)$, all estimated by counting in a single pass---$O(nd)$ training, streaming-friendly, embarrassingly parallel. \textbf{Always work in the log domain}: a product of a few hundred probabilities of order $10^{-4}$ underflows float64 (which bottoms out near $10^{-308}$) long before a document is scored, and $\log$-sum-exp recovers normalized posteriors when they are needed.
\end{mathresult}

The assumption is essentially always false---in text, ``San'' and ``Francisco'' are not conditionally independent given the topic---and the model works anyway, which is the interesting part.

\begin{keyinsight}[Right Argmax, Wrong Probabilities]
Classification only requires the \emph{ranking} of $\Pr(c\given x)$ to be correct, not its values. Dependent features amount to counting the same evidence multiple times, which drives the winning posterior exponentially toward $1$ while usually leaving the argmax alone (Domingos \& Pazzani, 1997). So Naive Bayes can be a decent classifier and a terrible probability estimator at the same time: its outputs pile up at $0$ and $1$ and are badly \textbf{miscalibrated}. Treating a Naive Bayes score as a probability---thresholding it for an expected-value decision, feeding it into an auction bid, showing it to a user---is a red flag. Calibrate it (isotonic or Platt) before it is allowed to mean anything numeric; see Chapter~\ref{cml:chap:probabilistic_foundations} and [DL 23].
\end{keyinsight}

\subsection{Smoothing and Variants}

\textbf{Laplace (add-$\alpha$) smoothing} is not an optional refinement---without it the model is broken. If a token never appears in class $c$ in training, its MLE estimate is $\Pr(x_j\given c) = 0$, and a single zero \emph{annihilates the entire product}: one unseen word vetoes fifty pieces of strong evidence. The fix is
\begin{equation}
\Pr(x_j = v\given c) = \frac{N_{v,c} + \alpha}{N_c + \alpha V},
\end{equation}
with $\alpha = 1$ the Laplace default and $\alpha < 1$ (Lidstone) common for large vocabularies. Concretely, with a $20{,}000$-word vocabulary and $100{,}000$ spam tokens, an unseen word moves from $0$ to $1/(100{,}000 + 20{,}000) = 8.3\times10^{-6}$: small enough not to distort the ranking, nonzero enough not to destroy it. Equivalently, add-$\alpha$ is the posterior mean under a symmetric Dirichlet$(\alpha)$ prior---the MAP framing developed in Chapter~\ref{cml:chap:probabilistic_foundations}.

\textbf{A worked score}, because interviewers ask for the arithmetic. Take priors $\Pr(\text{spam})=0.3$ and $\Pr(\text{ham})=0.7$, and a three-token document \{\emph{free}, \emph{meeting}, \emph{click}\} with likelihoods $(0.05,\,0.001,\,0.03)$ under spam and $(0.002,\,0.02,\,0.004)$ under ham. In logs:
\begin{align*}
\text{spam:}\quad & \ln 0.3 + \ln 0.05 + \ln 0.001 + \ln 0.03 = -1.204 - 2.996 - 6.908 - 3.507 = -14.614,\\
\text{ham:}\quad & \ln 0.7 + \ln 0.002 + \ln 0.02 + \ln 0.004 = -0.357 - 6.215 - 3.912 - 5.521 = -16.005 .
\end{align*}
Spam wins by $1.391$ nats, so the posterior odds are $e^{1.391} = 4.02$ and $\Pr(\text{spam}\given \text{doc}) = 4.02/5.02 = 0.80$. Now delete the smoothing: had \emph{meeting} never appeared in a spam training document, $\Pr(\text{meeting}\given\text{spam}) = 0$, the spam log-score is $-\infty$, and the document is classified ham no matter how damning the other two tokens are. One unseen token vetoes the entire model---which is the whole argument for add-$\alpha$ in a single sentence.

The three variants differ only in the per-feature likelihood:
\begin{itemize}
    \item \textbf{Multinomial NB}: features are counts; $\Pr(x\given c)$ is a multinomial over the vocabulary. The default for text with term frequencies. Its log-posterior is \emph{linear} in the count vector, so multinomial NB is a linear classifier whose weights happen to be estimated by counting instead of by optimization.
    \item \textbf{Bernoulli NB}: features are binary presence indicators, and the likelihood includes an explicit term for \emph{absent} features, $\prod_{j \notin x}(1 - \Pr(x_j\given c))$. Better on short documents; different from multinomial NB on the same binarized data precisely because of that absence term.
    \item \textbf{Gaussian NB}: continuous features, with a per-class per-feature $\mathcal{N}(\mu_{jc}, \sigma_{jc}^2)$. Note the boundary geometry: shared variances across classes give a \emph{linear} boundary; per-class variances give a \emph{quadratic} one---Gaussian NB is diagonal-covariance QDA.
\end{itemize}

\subsection{Generative Versus Discriminative}

Naive Bayes and logistic regression are the canonical pair, and the comparison is a standing interview question.

\begin{itemize}
    \item A \textbf{generative} model learns $\Pr(x, y)$---typically $\Pr(y)$ and $\Pr(x\given y)$---and derives $\Pr(y\given x)$ through Bayes' rule. It models how the data was produced.
    \item A \textbf{discriminative} model learns $\Pr(y\given x)$ (or just a decision boundary) directly, and never claims to know anything about how $x$ is distributed.
\end{itemize}

They are a matched pair in a precise sense: under Naive Bayes' own assumptions the induced $\Pr(y\given x)$ has \emph{exactly} the logistic form, so the two models share a hypothesis class and differ only in how the parameters are fitted---by counting under the generative assumption, or by maximizing conditional likelihood without it.

\begin{mathresult}{Ng \& Jordan (2001): the Sample-Complexity Trade}
The generative model has a \textbf{higher asymptotic error} (its independence assumption is a bias that more data cannot remove), but \textbf{converges to that asymptote much faster}: Naive Bayes needs $O(\log d)$ training examples to approach its asymptotic error, logistic regression needs $O(d)$. The learning curves therefore cross. The practical rule: \textbf{generative wins in the small-$n$, large-$d$ regime; discriminative wins once you have enough data}, and ``enough'' scales with the number of features.
\end{mathresult}

The rest of the trade-off, in the order worth delivering:
\begin{itemize}
    \item \textbf{Generative buys structure.} Missing features are handled by \emph{marginalizing them out} rather than imputing---a real operational advantage. The model can sample new $x$, supports semi-supervised learning from unlabeled data, and gives an outlier/novelty score through $\Pr(x)$. Class priors can be adjusted at inference without refitting (swap $\Pr(c)$ for the deployment base rate), which is a genuinely useful lever under label shift.
    \item \textbf{Discriminative buys accuracy per assumption.} It spends all its capacity on the boundary, tolerates arbitrarily correlated and redundant features, and does not waste parameters modeling aspects of $\Pr(x)$ irrelevant to the decision. Vapnik's dictum applies: do not solve a harder problem (estimating a full density) than the one you need.
    \item \textbf{The framing generalizes.} GDA/QDA versus logistic regression, HMM versus CRF, and---at the modern end---a language model that maximizes $\Pr(\text{tokens})$ versus a reward model or classifier trained on top of it. A candidate who can carry the axis beyond the NB/LR pair is doing more than reciting.
\end{itemize}

Where Naive Bayes still earns its place in 2026: as a five-line baseline that trains in one pass over data too big to fit in memory, in high-dimensional sparse text with very few labels, and as the fast first filter in a cascade. What it is not is a serious contender against a tuned GBDT or a fine-tuned encoder on any problem where you have the labels for one.

%==============================================================================
\section{$k$-Nearest Neighbors}
\label{cml:sec:knn}
%==============================================================================

$k$-NN is the purest \emph{nonparametric} method and the cleanest place in the classical canon to see the curse of dimensionality bite. Its cost structure is inverted relative to everything else: \textbf{training is free} (store the data) and \textbf{inference is expensive}---$O(nd)$ per query for brute-force search, with $O(nd)$ memory that never shrinks. Every practical concern follows from that inversion.

\subsection{$k$ Is the Bias--Variance Knob}

At $k=1$ the model interpolates: training error is exactly zero, the decision boundary is a noisy Voronoi tessellation, and a single mislabeled point creates its own region---maximum variance, minimum bias. At $k=n$ every query returns the global majority class (or the global mean), which is the class prior---maximum bias, zero variance. For regression with homoscedastic noise the variance of the prediction is $\sigma^2/k$ while the bias grows with the neighborhood radius, which makes the trade explicit; the decomposition itself is [DL 2]. Choose $k$ by cross-validation; use odd $k$ for binary problems to avoid ties.

The classical guarantee is worth knowing: as $n\to\infty$, the 1-NN error rate is at most \emph{twice} the Bayes error (Cover \& Hart, 1967). $k$-NN with $k\to\infty$ and $k/n\to 0$ is Bayes-consistent. Both results are asymptotic in $n$ \emph{at fixed $d$}, which is exactly the assumption the next subsection destroys.

\subsection{Why It Fails in High Dimensions}

\begin{mathresult}{Two Ways to Show the Curse}
\textbf{Volume concentrates in the shell.} The fraction of a $d$-dimensional unit ball's volume lying in the outer shell of thickness $\varepsilon$ is $1 - (1-\varepsilon)^d$. At $d = 100$ and $\varepsilon = 0.1$, that is $1 - 0.9^{100} = 1 - 2.7\times10^{-5} = 99.997\%$. Essentially all of the volume---and therefore essentially all of the data---sits at nearly the same distance from any query point.

\textbf{Neighborhoods stop being local.} To capture a fraction $r$ of uniformly distributed data in a $d$-dimensional unit hypercube, a cubical neighborhood must have edge length $r^{1/d}$. To capture $1\%$ of the data: edge $0.10$ in $d=2$, edge $0.63$ in $d=10$, edge $\mathbf{0.955}$ in $d=100$. In a hundred dimensions the ``local'' neighborhood spans $95\%$ of the range of every feature---it is not local, it is the whole dataset with extra steps.
\end{mathresult}

The consequence for the algorithm is \textbf{distance concentration}: for independent features, the ratio of the farthest to the nearest neighbor's distance tends to $1$ as $d$ grows (Beyer et al., 1999), so the ``nearest'' neighbor is barely nearer than a random point and the votes it casts are close to noise. Note the escape clause that a strong answer supplies: what matters is the \emph{intrinsic} dimension, not the ambient one. Data on a low-dimensional manifold inside a high-dimensional space is fine, which is precisely why $k$-NN over \emph{learned embeddings} works spectacularly well while $k$-NN over a thousand raw sparse features does not. $k$-NN did not die of the curse; it moved into representation space, where it is the retrieval half of every modern search and RAG stack [SR 3], [NLP 9].

\subsection{The Practicalities}

\begin{itemize}
    \item \textbf{Scaling is mandatory}, not advisory. Euclidean distance is dominated by whichever feature has the largest units; a salary in dollars and an age in years put the model entirely at the mercy of salary. Standardize (or min-max scale) inside the cross-validation pipeline.
    \item \textbf{Distance metric.} Euclidean by default; Manhattan is somewhat more robust in high dimensions; \textbf{cosine} for text and embeddings, where magnitude is uninformative; Hamming for categorical; Mahalanobis when features are correlated and you can estimate the covariance. Learning the metric instead of choosing it is [DL 7].
    \item \textbf{Weighting.} Distance-weighted voting (weight $\propto 1/\text{dist}$) makes $k$ far less sensitive and softens the boundary. Under class imbalance, plain majority voting is biased toward the majority class---weight the votes or resample.
    \item \textbf{Making queries fast.} Exact structures---kd-trees, ball trees---help only in low dimensions; a kd-tree degenerates to a linear scan past roughly $20$ dimensions because the pruning bound stops pruning. Beyond that you accept approximation: LSH, IVF, and HNSW graphs. The kd-tree construction and search are a coding-round staple (Chapter~\ref{cml:chap:coding_canon}); the production ANN indexes and their recall/latency knobs are [SR 3].
    \item \textbf{Where it is genuinely the right tool}: as a strong baseline on small tabular data, for on-the-fly personalization where the ``model'' must reflect data inserted seconds ago, in retrieval and recommendation over embeddings, and as an interpretable-by-example system (``here are the three most similar past cases''), which some regulated settings value highly.
\end{itemize}

%==============================================================================
\section{Choosing Among the Classical Models}
\label{cml:sec:classical_choice}
%==============================================================================

The models in this chapter are rarely the final answer on a modern problem---trees and neural networks usually are---but every one of them is still the right \emph{first} answer to some question, and the loop rewards knowing which. Table~\ref{cml:tab:model_card} is the compressed version; the ``start simple, add complexity'' framing it instantiates is [DL 24].

\begin{table}[H]
\centering\small
\caption{The classical model card. $n$ = examples, $d$ = features, $n_{\text{SV}}$ = support vectors, $k$ = optimizer iterations.}
\label{cml:tab:model_card}
\begin{tabularx}{\textwidth}{@{}l l l c c X@{}}
\toprule
\textbf{Model} & \textbf{Train} & \textbf{Predict} & \textbf{Scale?} & \textbf{Probs?} & \textbf{Still the right call when} \\
\midrule
OLS / ridge & $O(nd^2{+}d^3)$ & $O(d)$ & for ridge & --- & Small $d$, interpretable effects, collinearity (ridge). \\
Lasso / elastic net & $O(ndk)$ & $O(d_{\text{nz}})$ & yes & --- & Selection wanted; $d \gg n$; a small serving model. \\
Logistic regression & $O(ndk)$ & $O(d)$ & for penalty & calibrated & Interpretability or regulation; huge sparse online data; a baseline. \\
Linear SVM & $O(nd)$ per pass & $O(d)$ & yes & no (Platt) & Very high-dimensional sparse text, $d \gg n$. \\
Kernel SVM & $O(n^2)$--$O(n^3)$ & $O(n_{\text{SV}}d)$ & yes & no (Platt) & Small $n$, strong nonlinearity, a known good kernel. \\
Naive Bayes & $O(nd)$, one pass & $O(d)$ & no & poorly & Tiny labeled data, huge sparse text, streaming baseline. \\
$k$-NN & $O(1)$ & $O(nd)$ & yes & vote share & Low intrinsic dimension, embeddings, example-based explanation. \\
\bottomrule
\end{tabularx}
\end{table}

\subsection{The Five Whiteboard Derivations}

If a classical round goes deep, it goes deep on one of these five. Each should be producible on a whiteboard, unprompted, in under three minutes, with the geometry or intuition narrated while the algebra is written.

\begin{enumerate}
    \item \textbf{Normal equations} from $\nabla_\beta\norm{y-X\beta}^2 = 0$, with the orthogonality reading, the invertibility condition, and the closed-form-versus-gradient-descent cost comparison (Section~\ref{cml:sec:linear_regression}).
    \item \textbf{Logistic NLL and its gradient}: Bernoulli likelihood $\to$ cross-entropy $\to$ $X^{\top}(\sigma(X\beta)-y)$, with the $p(1-p)$ cancellation shown and convexity justified by $X^{\top}WX \succeq 0$ (Section~\ref{cml:sec:logistic_glm}).
    \item \textbf{L1 sparsity}, both halves: the corner geometry via the constrained form, and soft-thresholding via the subgradient in the orthonormal case (Section~\ref{cml:sec:regularization}).
    \item \textbf{SVM}: geometric margin under canonical scaling $\to$ primal QP $\to$ Lagrangian $\to$ dual in $\alpha$ $\to$ kernel trick, then identify support vectors from complementary slackness (Section~\ref{cml:sec:svm_kernels}).
    \item \textbf{Naive Bayes decision rule} from Bayes' rule plus conditional independence, in log space, with the smoothing argument attached (Section~\ref{cml:sec:naive_bayes}).
\end{enumerate}

\begin{keyinsight}[One Generator Behind Four of the Five]
Derivations 1, 2, and half of 3 are the same move: write the objective, take the gradient, set it to zero, and read off what the stationarity condition says. Derivation 4 is that move under constraints, which is what the Lagrangian is for. If a derivation is forgotten mid-interview, the recovery is to say so and rebuild from the objective---interviewers reward that visibly, because it is the clearest evidence that the knowledge is generative rather than memorized.
\end{keyinsight}

%==============================================================================
\section{Interview Questions}
\label{cml:sec:supervised_interview}
%==============================================================================

\begin{interviewq}
\textbf{Q: Why does L1 regularization produce sparse solutions and L2 does not?}

\badgeLfive{L5} \quad \badgetype{Mathematical} \quad \badgefreq{\freqhigh}

\stronganswer

Give both halves; either alone is a partial answer. \textbf{The geometry.} By Lagrangian duality the penalized problem $\min \tfrac12\norm{y-X\beta}^2 + \lambda P(\beta)$ is equivalent to minimizing RSS subject to $P(\beta) \le t$, so the estimate is the point where the smallest RSS ellipse first touches the constraint set. The L1 ball is a cross-polytope whose \emph{corners lie on the coordinate axes}; a corner accepts a whole cone of contour orientations, so contact there is the typical case, not a coincidence---and contact at a corner means a coefficient is exactly zero. The L2 ball is smooth, so the contact point almost surely has every coordinate nonzero: shrinkage without selection. \textbf{The algebra.} With an orthonormal design the problem separates and both have closed forms: ridge gives $\hat\beta_j^{\text{OLS}}/(1+\lambda)$, a uniform multiplicative shrink, while the lasso gives soft-thresholding, $\operatorname{sign}(\hat\beta_j^{\text{OLS}})(|\hat\beta_j^{\text{OLS}}|-\lambda)_+$. The mechanism is visible in the subgradient: at $\beta_j=0$ the subdifferential of $|\beta_j|$ is the whole interval $[-\lambda,\lambda]$, so zero is optimal whenever the correlation $|x_j^{\top}r|$ falls inside it. Ridge's derivative $\lambda\beta_j$ \emph{vanishes} as $\beta_j\to 0$, so it can never finish the job. Worked: $\hat\beta^{\text{OLS}}=(2.0,0.4,-0.9)$, $\lambda=0.5$ gives lasso $(1.5,\mathbf{0},-0.4)$ and ridge $(1.33,0.27,-0.60)$. Close with the Bayesian reading---L2 is a Gaussian-prior MAP, L1 a Laplace-prior MAP, and the Laplace density's spike at the origin is the corner in probabilistic clothing---and note that this same soft-threshold operator \emph{is} the coordinate-descent update that real lasso solvers run.

\redflags
\begin{itemize}
    \item ``L1 is a stronger penalty.'' It is not about strength; for small coefficients L1 is a \emph{weaker} penalty than L2 in absolute value. Sparsity comes from the non-differentiability at zero.
    \item Drawing the diamond and stopping. Interviewers at this level expect the soft-thresholding or subgradient argument as well.
    \item Claiming L2 ``sets small weights to zero eventually''---it approaches zero asymptotically and reaches it only at $\lambda=\infty$.
    \item Forgetting that the features must be standardized for either penalty to be meaningful.
\end{itemize}

\interviewtest{Whether regularization is a picture you memorized or a piece of mathematics you can operate. This is the most-asked classical question after trees, so a fluent, two-sided answer sets the tone for everything that follows; a hesitant one invites twenty minutes of drilling.}

\goodgreat
\begin{itemize}
    \item \badgeLfive{L5} The corner geometry stated correctly, plus the constrained/penalized equivalence and the fact that L2 shrinks but never zeroes.
    \item \badgeLsix{L6} Adds soft-thresholding with the subgradient argument, the numeric example, the Laplace/Gaussian MAP reading, and knows the lasso is biased for large coefficients.
    \item \badgeLseven{L7} Discusses when sparsity is the wrong goal (correlated groups, where the lasso's choice is arbitrary and unstable---hence elastic net), the $d>n$ selection cap, why post-lasso inference is invalid without a correction, and how the same machinery appears as structured sparsity or pruning in deep models.
\end{itemize}

\followups{``Derive the soft-threshold operator from the subgradient condition.'' ``Two features are correlated at 0.99---what does the lasso do, and why is that a problem?'' ``How would you choose $\lambda$, and what is the one-standard-error rule?''}
\end{interviewq}

\begin{interviewq}
\textbf{Q: Derive the normal equations for linear regression. When would you not use the closed form?}

\badgeLfive{L5} \quad \badgetype{Mathematical} \quad \badgefreq{\freqhigh}

\stronganswer

\textbf{Derivation.} Minimize $J(\beta)=\norm{y-X\beta}^2 = y^{\top}y - 2\beta^{\top}X^{\top}y + \beta^{\top}X^{\top}X\beta$. Then $\nabla_\beta J = -2X^{\top}y + 2X^{\top}X\beta$, and setting it to zero gives $X^{\top}X\hat\beta = X^{\top}y$, so $\hat\beta = (X^{\top}X)^{-1}X^{\top}y$ when $X$ has full column rank. The Hessian $2X^{\top}X$ is PSD, so the stationary point is a global minimum. Say the geometry as you write it: the condition is $X^{\top}(y-X\hat\beta)=0$, i.e., the residual is orthogonal to the column space---OLS is the orthogonal projection of $y$ onto $\mathrm{col}(X)$, with $H = X(X^{\top}X)^{-1}X^{\top}$ the projection matrix. \textbf{When not to.} Four separate reasons, not one: (i) \emph{cost}---$O(nd^2 + d^3)$ time and $O(d^2)$ memory, so at $d=10^4$ the Gram matrix alone is $800$\,MB and the multiply is $10^{14}$ flops, versus $O(ndk)$ for gradient descent; (ii) \emph{singularity}---if $d>n$ or any feature is an exact combination of others (the dummy-variable trap), $X^{\top}X$ has no inverse; (iii) \emph{conditioning}---$\kappa(X^{\top}X)=\kappa(X)^2$, so forming the Gram matrix costs you half your significant digits under near-collinearity; (iv) \emph{the setting}---streaming or distributed data, or a non-quadratic loss, admits no closed form at all. And the implementation point that separates candidates: \textbf{never invert}. Solve by Cholesky when the problem is well conditioned and $d\ll n$, by QR on $X$ directly for numerical safety (it works at $\kappa(X)$, not $\kappa(X)^2$), by SVD when $X$ may be rank-deficient---which returns the minimum-norm solution, and is exactly why \texttt{numpy.linalg.lstsq} returns an answer where the textbook formula would divide by zero. Ridge is the other fix: $X^{\top}X + \lambda I$ is positive definite for any $X$.

\redflags
\begin{itemize}
    \item Writing \texttt{inv(X.T @ X) @ X.T @ y} as the recommended implementation.
    \item Naming only ``it is slow'' and missing singularity and conditioning.
    \item Not knowing the geometric interpretation, or that the residual is orthogonal to the fitted values.
    \item Claiming OLS needs normally distributed features or targets---normality is about the \emph{errors}, and only for inference.
\end{itemize}

\interviewtest{Whether the most basic derivation in machine learning is genuinely at your fingertips, and whether you know the numerical-linear-algebra layer underneath it. The follow-up about implementation is where staff-level candidates separate from ones who have only used \texttt{fit()}.}

\goodgreat
\begin{itemize}
    \item \badgeLfive{L5} Correct derivation, the invertibility condition, and cost versus gradient descent.
    \item \badgeLsix{L6} All four failure modes distinguished, the QR/Cholesky/SVD ladder with the $\kappa^2$ argument, and the projection geometry with leverage.
    \item \badgeLseven{L7} Connects to ridge as the conditioning repair, notes the minimum-norm interpolator in the $d>n$ regime and its link to double descent [DL 2], and can size the choice for a concrete $(n,d)$ and a memory budget.
\end{itemize}

\followups{``What is $\kappa(X^{\top}X)$ versus $\kappa(X)$, and why does it matter?'' ``What does the SVD give you when $X$ is rank-deficient?'' ``How does leverage $h_{ii}$ let you compute LOOCV in one fit?''}
\end{interviewq}

\begin{interviewq}
\textbf{Q: Derive logistic regression's loss from maximum likelihood, give its gradient, and explain why we do not just use squared error.}

\badgeLfive{L5} \quad \badgetype{First Principles} \quad \badgefreq{\freqhigh}

\stronganswer

\textbf{Derivation.} Assume $y_i\given x_i \sim \mathrm{Bernoulli}(p_i)$ and model the log-odds as linear, $\log\frac{p_i}{1-p_i}=x_i^{\top}\beta$, so $p_i=\sigma(x_i^{\top}\beta)$---the sigmoid is the \emph{inverse link}, not an arbitrary squashing function. The likelihood $\prod_i p_i^{y_i}(1-p_i)^{1-y_i}$ has negative log $-\sum_i[y_i\log p_i + (1-y_i)\log(1-p_i)]$, which \emph{is} binary cross-entropy: log loss is the Bernoulli NLL, not a separate choice. \textbf{Gradient.} $\partial\loss_i/\partial p_i = (p_i-y_i)/(p_i(1-p_i))$ and $\partial p_i/\partial z_i = p_i(1-p_i)$; the factors cancel, leaving $\partial\loss_i/\partial z_i = p_i - y_i$ and $\nabla_\beta\loss = X^{\top}(\sigma(X\beta)-y)$---identical in shape to least squares' $X^{\top}(X\beta - y)$, which is the GLM canonical-link result, not a coincidence. The Hessian is $X^{\top}WX$ with $W=\diag(p_i(1-p_i))\succeq 0$, so the problem is convex with no bad local minima; there is no closed form, and the classical fit is Newton/IRLS, with L-BFGS the modern default and SGD/FTRL at scale. \textbf{Why not squared error.} Three separate reasons. (1) Likelihood: squared error is the Gaussian NLL and the label is Bernoulli. (2) Convexity: $\tfrac12(\sigma(z)-y)^2$ is not convex in $\beta$---with one feature and $y=0$, $f(\beta)=\tfrac12\sigma(\beta)^2$ has $f''=\sigma^2(1-\sigma)(2-3\sigma)<0$ once $\sigma>2/3$. (3) Gradient saturation: the squared-error gradient carries an extra $\sigma'(z)$ factor, so at $z=-10$ with $y=1$---confidently wrong---it is about $4.5\times10^{-5}$ against log loss's $-0.99995$, roughly $22{,}000\times$ smaller. The example that most deserves a big update contributes almost nothing. Add the nuance: squared error on probabilities is the \emph{Brier score}, which \emph{is} a proper scoring rule, so the objection is about likelihood and optimization geometry, not calibration propriety.

\redflags
\begin{itemize}
    \item ``Logistic regression has a closed-form solution.''
    \item Starting from ``we apply a sigmoid'' with no likelihood, then being unable to say where the loss came from.
    \item ``MSE gives uncalibrated probabilities''---the Brier score is proper; this is a confidently wrong claim.
    \item Writing the gradient with a stray $\sigma'$ factor, which signals the cancellation was memorized rather than derived.
\end{itemize}

\interviewtest{The canonical whiteboard derivation of the classical loop. Interviewers watch for the cancellation step, the convexity claim with a reason attached, and whether ``why not MSE'' produces one vague sentence or three sharp ones.}

\goodgreat
\begin{itemize}
    \item \badgeLfive{L5} Bernoulli $\to$ NLL $\to$ cross-entropy, the clean gradient, convexity, no closed form.
    \item \badgeLsix{L6} Adds the Hessian and IRLS, the non-convexity counterexample, the saturation arithmetic, and the Brier nuance.
    \item \badgeLseven{L7} Places it in the GLM frame---canonical link gives $b'(\theta)=\mu$ and hence the $(\mu - y)$ gradient---and connects to softmax cross-entropy [DL 1] and to gradient boosting's log-loss pseudo-residuals $y_i - p_i$ (Chapter~\ref{cml:chap:trees_ensembles}).
\end{itemize}

\followups{``Show the Hessian and prove convexity.'' ``What is IRLS actually doing?'' ``How does this change for $K$ classes, and what is the identifiability issue?''}
\end{interviewq}

\begin{interviewq}
\textbf{Q: Derive the SVM: margin, primal, dual, and the kernel trick. What exactly is a support vector?}

\badgeLsix{L6} \quad \badgetype{Mathematical} \quad \badgefreq{\freqmed}

\stronganswer

\textbf{Margin.} With $f(x)=w^{\top}x+b$ and $y\in\{-1,+1\}$, the distance from $x_i$ to the hyperplane is $y_i f(x_i)/\norm{w}$. Since $(w,b)$ can be rescaled freely, impose the canonical normalization $\min_i y_i f(x_i)=1$; then the margin is $1/\norm{w}$, the band is $2/\norm{w}$ wide, and maximizing it is $\min \tfrac12\norm{w}^2$ subject to $y_i f(x_i)\ge 1$---a convex QP. \textbf{Soft margin.} Add slack: $\min\tfrac12\norm{w}^2 + C\sum\xi_i$ with $y_i f(x_i)\ge 1-\xi_i,\ \xi_i\ge 0$. At the optimum $\xi_i=\max(0,1-y_if(x_i))$, so eliminating the constraints gives $\min\tfrac12\norm{w}^2 + C\sum\max(0,1-y_if(x_i))$: hinge loss with an L2 penalty, $\lambda=1/(2C)$. Being able to move between the constrained and loss forms is the point. \textbf{Dual.} Build the Lagrangian with $\alpha_i\ge0$ on the margin constraints and $\mu_i\ge0$ on the slacks. Stationarity gives $w=\sum_i\alpha_iy_ix_i$ (the representer form), $\sum_i\alpha_iy_i=0$, and $\alpha_i+\mu_i=C$, which becomes the box $0\le\alpha_i\le C$. Substituting yields $\max_\alpha \sum_i\alpha_i - \tfrac12\sum_{i,j}\alpha_i\alpha_jy_iy_j\langle x_i,x_j\rangle$ under those constraints, and $f(x)=\sum_i\alpha_iy_i\langle x_i,x\rangle+b$. \textbf{Kernel trick.} The data now appears \emph{only} through inner products, so replace $\langle x_i,x_j\rangle$ with any PSD kernel $K(x_i,x_j)=\langle\phi(x_i),\phi(x_j)\rangle$ and you have fitted a linear model in $\phi$-space without ever computing $\phi$---for the RBF kernel you could not, since expanding $e^{2\gamma x^{\top}x'}=\sum_k (2\gamma)^k(x^{\top}x')^k/k!$ shows the feature space is infinite-dimensional. \textbf{Support vectors, from KKT.} Complementary slackness $\alpha_i[y_if(x_i)-1+\xi_i]=0$ gives three cases: $\alpha_i=0$ means the point is strictly outside the margin and irrelevant (delete it, nothing changes); $0<\alpha_i<C$ means $\xi_i=0$ and $y_if(x_i)=1$, a free support vector sitting exactly on the margin, which is what you use to recover $b$; $\alpha_i=C$ means the point is inside the margin or misclassified. Support vectors are the points with $\alpha_i>0$---sparsity \emph{in examples}, which is what makes inference $O(n_{\text{SV}}d)$.

\redflags
\begin{itemize}
    \item ``The kernel trick computes high-dimensional features efficiently''---it never computes them at all.
    \item Cannot state what a support vector is, or defines it as ``the closest points'' without the $\alpha_i>0$ / KKT characterization (which also covers margin violators).
    \item ``SVMs maximize accuracy.'' They maximize a regularized margin; hinge loss is a convex surrogate for 0--1 loss.
    \item Getting $C$ backwards (thinking large $C$ means more regularization).
    \item Not knowing any condition for a valid kernel, or asserting any similarity function will do.
\end{itemize}

\interviewtest{Optimization literacy: constrained versus penalized formulations, Lagrangian duality, and reading structure off KKT conditions---in the one setting small enough for a whiteboard. Interviewers rarely expect a production SVM; they expect this derivation.}

\goodgreat
\begin{itemize}
    \item \badgeLfive{L5} Margin $=1/\norm{w}$ under canonical scaling, the primal QP, and the kernel trick stated correctly.
    \item \badgeLsix{L6} The full dual with the box constraint derived from $\alpha_i+\mu_i=C$, the three KKT cases, the hinge-loss equivalence, and Mercer's PSD condition.
    \item \badgeLseven{L7} Volunteers the honest 2026 position---$O(n^2)$ kernel matrix, no native probabilities, GBDTs and NNs won---while noting where kernel thinking survives: Gaussian processes, the kernel-regression reading of attention [NLP 5], NTK theory, and random-feature approximations as the scaling escape.
\end{itemize}

\followups{``Why is the dual bounded by $C$ and not just $\alpha_i \ge 0$?'' ``How do you recover $b$?'' ``When would you solve the primal instead of the dual?''}
\end{interviewq}

\begin{interviewq}
\textbf{Q: Generative versus discriminative---define them, and use Naive Bayes and logistic regression to explain the trade-off.}

\badgeLsix{L6} \quad \badgetype{Conceptual} \quad \badgefreq{\freqmed}

\stronganswer

\textbf{Definitions.} A generative model learns the joint $\Pr(x,y)$---in practice $\Pr(y)$ and $\Pr(x\given y)$---and inverts it with Bayes' rule; a discriminative model learns $\Pr(y\given x)$ directly and makes no claim about how $x$ is distributed. \textbf{Why NB and LR are the canonical pair}: under Naive Bayes' conditional-independence assumption, the induced posterior $\Pr(y\given x)$ has \emph{exactly} the logistic form. Same hypothesis class, different fitting principle---counting under a generative assumption versus maximizing conditional likelihood without one. \textbf{The trade-off} (Ng \& Jordan, 2001): Naive Bayes has the higher asymptotic error, because the independence assumption is a bias that data cannot remove, but it reaches its asymptote in $O(\log d)$ examples where logistic regression needs $O(d)$. The learning curves cross: generative wins in the small-$n$, large-$d$ regime and discriminative wins once labels are plentiful. \textbf{What else generative buys}: missing features handled by \emph{marginalization} instead of imputation; the ability to sample data; a density $\Pr(x)$ usable for novelty detection; semi-supervised use of unlabeled data; and class priors that can be swapped at inference time to match a new base rate without refitting---a genuinely useful lever under label shift. \textbf{What discriminative buys}: it spends all capacity on the boundary, tolerates correlated and redundant features that would double-count under NB, and does not waste parameters modeling irrelevant aspects of $\Pr(x)$---Vapnik's ``do not solve a harder problem than you need.'' Close by generalizing the axis: GDA versus logistic regression, HMM versus CRF, and a language model maximizing $\Pr(\text{tokens})$ versus a classifier or reward model trained on top of it. And add the practical caveat: NB's posteriors are badly miscalibrated because dependent features double-count evidence, so its scores need calibration before they are allowed to mean anything numeric.

\redflags
\begin{itemize}
    \item ``Generative models generate data, discriminative models discriminate''---a restatement, not a definition. Say $\Pr(x,y)$ versus $\Pr(y\given x)$.
    \item Claiming Naive Bayes is simply worse. The whole point is that its bias buys sample efficiency, and the curves cross.
    \item Treating Naive Bayes probabilities as calibrated.
    \item Not knowing that NB and LR share a functional form for the posterior, which is what makes them \emph{the} comparison.
\end{itemize}

\interviewtest{Whether you hold a framework or two model descriptions. The Ng \& Jordan result and the missing-data/label-shift advantages are the markers that the axis is understood rather than recited.}

\goodgreat
\begin{itemize}
    \item \badgeLfive{L5} Correct definitions, the NB/LR pairing, and ``NB for small data, LR for large.''
    \item \badgeLsix{L6} Adds the asymptotic-error-versus-convergence-rate statement with the $\log d$ / $d$ scaling, marginalization over missing features, and the calibration caveat.
    \item \badgeLseven{L7} Extends the axis to sequence models and to modern generative pretraining plus discriminative heads, and discusses when a hybrid is right (generative pretraining for representation, discriminative fine-tuning for the decision).
\end{itemize}

\followups{``Naive Bayes' independence assumption is clearly false for text---why does it still work?'' ``Which model would you pick with 200 labeled documents and 50{,}000 features?'' ``How would you exploit the generative model's ability to handle missing features?''}
\end{interviewq}

\begin{interviewq}
\textbf{Q: Your logistic regression's weights blew up to $\pm 10^6$ and training accuracy is exactly 1.0. What happened, and what do you do?}

\badgeLfive{L5} \quad \badgetype{Debugging} \quad \badgefreq{\freqmed}

\stronganswer

\textbf{Diagnosis: complete (or quasi-complete) separation.} Some hyperplane perfectly splits the classes, so for any separating $\beta$ the scaled vector $c\beta$ pushes every fitted probability toward the correct $0$ or $1$ and the likelihood increases monotonically in $c$---\textbf{the MLE does not exist}, and the optimizer marches toward infinity until it hits an iteration cap. Quasi-complete separation, usually caused by a rare category that only ever appears with one label, produces the same divergence in a subset of coordinates. Confirming symptoms: enormous standard errors alongside the enormous coefficients, fitted probabilities pinned at $0$ and $1$, convergence warnings, and coefficients that keep growing if you raise \texttt{max\_iter}. \textbf{Response, in order.} (1) \emph{Check for leakage first.} A feature that perfectly predicts the label is more often a post-outcome variable---a field populated after the event, an ID correlated with the label pipeline, a duplicated target---than a discovery. Look at feature timestamps before celebrating. (2) \emph{Regularize.} Any L2 penalty makes the objective strongly convex, so the optimum is finite and unique; this is why scikit-learn regularizes by default, and turning the penalty off is usually how people meet this bug at all. (3) \emph{Firth's penalized likelihood} (a Jeffreys prior) is the statistician's fix and also removes small-sample bias, which matters if you need interpretable coefficients and standard errors rather than just predictions. (4) \emph{Handle the offending category}---collapse rare levels, add a smoothed prior---rather than blindly dropping the feature. Finally, note what does \emph{not} fix it: more iterations, a smaller learning rate, or a different optimizer. This is not an optimization failure; the objective genuinely has no finite minimizer.

\redflags
\begin{itemize}
    \item ``The learning rate is too high'' or ``it needs more epochs''---misreading a nonexistent optimum as an optimization bug.
    \item Not knowing the term \emph{separation}, or unable to say why the likelihood is maximized at infinity.
    \item Jumping to ``add L2'' without checking whether the perfect separator is a leaked feature. Regularizing a leak produces a well-conditioned model that is still wrong.
    \item Concluding the model is excellent because training accuracy is $1.0$.
\end{itemize}

\interviewtest{Whether you can reason from an optimization symptom back to a property of the objective, and whether your first instinct on suspiciously perfect performance is suspicion. The leakage branch is the seniority marker.}

\goodgreat
\begin{itemize}
    \item \badgeLfive{L5} Names separation, explains the infinite-likelihood argument, proposes regularization.
    \item \badgeLsix{L6} Distinguishes complete from quasi-complete separation, checks leakage first, knows Firth's method and why standard errors are also broken.
    \item \badgeLseven{L7} Discusses what separation implies for deployment (a model whose scores are all $0$ or $1$ is unusable for ranked or thresholded decisions), how to detect it automatically in a retraining pipeline, and the connection to margin maximization---L2-regularized logistic regression on separable data converges in direction to the max-margin solution as $\lambda\to0$.
\end{itemize}

\followups{``What if only one coefficient diverges?'' ``How would you detect separation before fitting?'' ``Would ridge or lasso be better here, and why?''}
\end{interviewq}

\begin{interviewq}
\textbf{Q: What does $C$ do in an SVM? What about $\gamma$ in an RBF kernel?}

\badgeLfive{L5} \quad \badgetype{Conceptual} \quad \badgefreq{\freqmed}

\stronganswer

Both are bias--variance dials, but on different axes, and the strong answer says which. \textbf{$C$ is inverse regularization.} In $\min\tfrac12\norm{w}^2 + C\sum\xi_i$, $C$ prices margin violations. Large $C$ makes slack expensive, driving toward a hard margin that bends to classify every training point---low bias, high variance, and acute sensitivity to mislabeled points near the boundary. Small $C$ tolerates violations for a wider, smoother margin---high bias, low variance. The hinge-loss form makes it exact: dividing by $C$ gives regularization strength $\lambda = 1/(2C)$, so $C$ is literally $1/(2\lambda)$. \textbf{$\gamma$ is the inverse kernel width.} With $K(x,x')=\exp(-\gamma\norm{x-x'}^2)$ and $\gamma=1/(2\sigma^2)$, large $\gamma$ makes each training point's influence a narrow spike---the boundary becomes wiggly and, in the limit, every point becomes its own support vector and the model memorizes the training set at chance-level generalization. Small $\gamma$ makes the kernel nearly constant over the data, so the model degenerates toward linear (high bias). \textbf{They interact}, which is why they are tuned jointly on a log-spaced grid rather than one at a time: a large $\gamma$ with a small $C$ can look acceptable because the regularization damps the wiggliness, and the pathologies show up only in the corners of the grid. Anchor with the default $\gamma = 1/(d\cdot\Var(X))$ (scikit-learn's \texttt{'scale'}) and a grid such as $C\in\{10^{-2},\dots,10^{3}\}$, $\gamma\in\{10^{-4},\dots,10^{1}\}$. \textbf{The precondition nobody should need reminding of}: features must be standardized. RBF distance is dominated by whatever feature has the largest units, so an unscaled SVM is effectively a one-feature model and no amount of grid search rescues it.

\redflags
\begin{itemize}
    \item Getting $C$ backwards---saying large $C$ means more regularization.
    \item ``$\gamma$ controls the learning rate'' or any answer that does not connect it to the kernel's width.
    \item Tuning them independently, or not mentioning feature scaling.
    \item Unable to say which direction of each knob overfits.
\end{itemize}

\interviewtest{A fast fluency check on the two hyperparameters everyone claims to have tuned. Reversing $C$ is common and immediately costly, because it suggests the hinge-loss form was never internalized.}

\goodgreat
\begin{itemize}
    \item \badgeLfive{L5} Both directions correct, both framed as bias--variance dials.
    \item \badgeLsix{L6} Derives $\lambda=1/(2C)$ from the hinge form, explains the $\gamma\to\infty$ memorization limit in terms of every point becoming a support vector, and insists on scaling and joint tuning.
    \item \badgeLseven{L7} Discusses the diagnostic value of the support-vector count (a model where nearly every point is a support vector is overfitting and will also be slow at inference), and the compute budget of a 2-D grid search versus randomized or Bayesian search on a large dataset.
\end{itemize}

\followups{``How would you recognize overfitting from the fitted model alone, without a validation set?'' ``What happens as $\gamma\to 0$?'' ``How does the support-vector count affect serving latency?''}
\end{interviewq}

\begin{interviewq}
\textbf{Q: Why does $k$-NN degrade in high dimensions? Quantify it.}

\badgeLfive{L5} \quad \badgetype{Conceptual} \quad \badgefreq{\freqmed}

\stronganswer

Because in high dimensions ``nearest'' stops meaning anything. Two quantifications, either of which should be produced on demand. \textbf{Volume concentrates in the shell}: the fraction of a $d$-dimensional unit ball's volume within the outer shell of thickness $\varepsilon$ is $1-(1-\varepsilon)^d$; at $d=100,\ \varepsilon=0.1$ that is $1-0.9^{100}\approx 99.997\%$. Essentially all the data sits at nearly the same distance from any query. \textbf{Neighborhoods stop being local}: to capture a fraction $r$ of uniform data in the unit hypercube, a cubical neighborhood needs edge length $r^{1/d}$---for $1\%$ of the data that is edge $0.10$ at $d=2$, $0.63$ at $d=10$, and $0.955$ at $d=100$. Your ``local'' neighborhood spans $95\%$ of every feature's range, so the prediction is barely different from the global average. The formal statement is \textbf{distance concentration} (Beyer et al., 1999): for independent features the ratio of farthest to nearest neighbor distance tends to $1$, so neighbor identity is dominated by noise and the votes are close to random. Add the cost problem---$O(nd)$ per query with no training-time amortization---and exact indexes do not save you: a kd-tree degenerates to a linear scan past roughly $20$ dimensions because its pruning bound stops pruning. \textbf{The modern resolution, which is the part that separates answers}: what matters is the \emph{intrinsic} dimension, not the ambient one. Data concentrated on a low-dimensional manifold is fine, which is exactly why $k$-NN over learned embeddings works spectacularly---a $768$-dimensional sentence embedding has an intrinsic dimension far below $768$, and the metric was trained so that neighbors mean something. $k$-NN did not die of the curse; it moved into representation space, where it is the retrieval half of every search and RAG stack [SR 3], [NLP 9]. The practical toolkit is therefore: reduce dimension (PCA, feature selection, Chapter~\ref{cml:chap:unsupervised}), \emph{learn} the metric [DL 7], and use approximate indexes (LSH, IVF, HNSW) rather than exact search.

\redflags
\begin{itemize}
    \item ``It gets slow''---true but secondary; the statistical failure is the point.
    \item Naming ``the curse of dimensionality'' with no mechanism and no number.
    \item Believing more data fixes it. Beating the curse needs exponentially more data in $d$, which is exactly why it is a curse.
    \item Not knowing that modern embedding-based retrieval is $k$-NN, which suggests the concept is filed as history.
\end{itemize}

\interviewtest{Whether you can quantify a piece of folklore. Any candidate can say ``curse of dimensionality''; the shell-volume or edge-length arithmetic, delivered in thirty seconds, is what distinguishes understanding, and the embedding update shows the knowledge is current.}

\goodgreat
\begin{itemize}
    \item \badgeLfive{L5} Distance concentration explained plus one quantitative argument, and feature scaling named as a precondition.
    \item \badgeLsix{L6} Both quantifications, the intrinsic-versus-ambient distinction, the kd-tree breakdown threshold, and embeddings as the modern fix.
    \item \badgeLseven{L7} Discusses the ANN recall/latency trade in production, when brute force on a GPU beats an index, and how metric learning changes the geometry rather than merely the dimension.
\end{itemize}

\followups{``What is the intrinsic dimension of your data and how would you estimate it?'' ``When is brute-force search actually the right choice?'' ``How does $k$ interact with dimensionality?''}
\end{interviewq}

\begin{interviewq}
\textbf{Q: When would you deploy logistic regression instead of gradient boosting in 2026?}

\badgeLsix{L6} \quad \badgetype{Trade-off} \quad \badgefreq{\freqhigh}

\stronganswer

Refuse the implied ``never,'' and give the conditions. \textbf{(1) Regulated or contested decisions.} Credit, insurance, hiring, and clinical settings often require a model whose behavior can be stated as a monotone effect per feature and audited by people who are not machine-learning engineers; a coefficient table with odds ratios is an artifact a regulator accepts, and adverse-action reasons fall out of it directly. Scorecards (binned features plus logistic regression) exist because of this, not because nobody tried boosting. \textbf{(2) Small data.} At a few thousand rows with a genuinely near-linear signal, a regularized linear model often matches a tuned GBDT and is far more stable across refits---and the variance of the \emph{model selection procedure} matters as much as the model. \textbf{(3) Extreme-scale sparse online serving.} The ad-CTR heritage: hundreds of millions of hashed sparse features, billions of daily examples, and hourly or continuous retraining. Logistic regression with FTRL-Proximal trains online in one pass, gives L1-induced sparsity that keeps the served model small, updates incrementally as data arrives, and scores in microseconds. GBDTs do not do online updates well and do not love ultra-high-cardinality sparse inputs. \textbf{(4) Well-calibrated probabilities out of the box.} Log loss is the Bernoulli NLL, so a fitted logistic model is usually close to calibrated without a second stage---valuable when the score is multiplied by a value to make an expected-value decision (bidding, expected loss). \textbf{(5) As the baseline that sets the bar.} It is the rung a staff engineer insists on before approving a complex model, and it makes the complexity budget explicit [DL 24]. \textbf{Then be honest about the default}: on mid-scale tabular data with nonlinearities and interactions, GBDTs win, usually by a lot, and the reasons are inductive-bias reasons (Chapter~\ref{cml:chap:trees_ensembles}). The good answer also notes the hybrid moves---feature crosses and splines to give the linear model curvature, or a GBDT used for feature discovery with a linear model in serving (the GBDT-plus-LR pattern), which buys some of the nonlinearity while keeping the serving path simple.

\redflags
\begin{itemize}
    \item ``Never, boosting is strictly better''---the failure the question is designed to catch.
    \item ``Logistic regression is more interpretable'' as the entire answer, with no mention of when interpretability is actually a requirement or what artifact it produces.
    \item Not knowing the online/sparse regime at all, which is where linear models still run at the largest scale in the industry.
    \item Claiming GBDT probabilities are uncalibrated by nature (they are often close, and calibration is measurable and fixable either way).
\end{itemize}

\interviewtest{Judgment rather than fashion: whether you can name concrete regimes where the simpler model is correct, and whether you concede the default honestly instead of contrarian posturing.}

\goodgreat
\begin{itemize}
    \item \badgeLfive{L5} Interpretability, small data, and speed, with the GBDT default conceded.
    \item \badgeLsix{L6} Adds the extreme-scale online sparse regime with FTRL and L1-for-serving-size, the calibration argument, and the hybrid patterns.
    \item \badgeLseven{L7} Frames it as a total-cost decision---retraining cadence, monitoring surface, on-call burden, model governance---and proposes the empirical test (fit both with an honest temporal split; if the gap is under the noise band, ship the simpler one).
\end{itemize}

\followups{``What is FTRL-Proximal doing that plain SGD is not?'' ``How would you give a linear model nonlinearity without losing interpretability?'' ``How do you show a regulator what the model does?''}
\end{interviewq}

\begin{interviewq}
\textbf{Q: You fit a linear model on 150 rows and 200 features. The coefficients flip sign between refits and two features have VIF over 40. Walk me through it.}

\badgeLsix{L6} \quad \badgetype{Debugging} \quad \badgefreq{\freqmed}

\stronganswer

Separate two distinct problems that happen to co-occur. \textbf{Problem 1: $d > n$, so OLS is not identified at all.} $X^{\top}X$ is singular, the minimizer is an affine set rather than a point, and whatever number came back is an artifact of the solver---\texttt{lstsq} returns the minimum-norm solution via SVD, which is \emph{a} least-squares solution but not \emph{the} answer to any question you asked. Training $R^2$ will be near $1$ and meaningless: with $d>n$ you can interpolate. \textbf{Problem 2: near-collinearity}, which would destabilize coefficients even at $n>d$. $\text{VIF}_j = 1/(1-R_j^2)$, so VIF $=40$ means $R_j^2 = 0.975$---that feature is $97.5\%$ explained by the others, and its coefficient's standard error is inflated $\sqrt{40}\approx 6.3\times$. Sign flips across bootstrap resamples are the expected consequence, and they are a statement about identifiability, not about the world. \textbf{The response.} (1) Ask what the coefficients are \emph{for}. If the goal is prediction, collinearity is largely harmless and the fix is regularization plus honest validation. If the goal is inference about individual effects, then with $150$ rows and $200$ correlated features that question is not answerable from this data, and saying so is the senior move. (2) Fit \textbf{ridge}: $X^{\top}X+\lambda I$ is positive definite for any $X$, so the solution exists and is unique even at $d>n$, and it directly repairs the conditioning; choose $\lambda$ by cross-validation. (3) If a small feature set is needed, use \textbf{elastic net} rather than lasso---with correlated groups the lasso picks one member essentially at random, so its selections are unstable across resamples, while the ridge component gives the grouping effect; at $d>n$ the lasso also caps out at $n$ selected features. (4) Report \textbf{stability}: run the selection under bootstrap or repeated CV and show selection frequencies (stability selection) rather than a single feature list. (5) Validate honestly---nested or repeated CV, since with $n=150$ a single split has enormous variance and hyperparameter selection on the same fold is itself an overfitting channel (Chapter~\ref{cml:chap:applied_craft}). (6) Consider reducing dimension first (PCA, or domain-driven grouping) or simply collecting more rows, which is the only thing that changes the underlying answerability.

\redflags
\begin{itemize}
    \item Reporting the OLS coefficients and interpreting them without noticing $d>n$.
    \item ``Drop the high-VIF features''---a defensible tactic sometimes, but not a diagnosis, and it silently changes what every remaining coefficient means.
    \item Using the lasso and presenting its selected set as \emph{the} important features, with no stability check.
    \item Trusting a single train/test split at $n=150$.
\end{itemize}

\interviewtest{Whether you distinguish identifiability from prediction, and whether you will tell a stakeholder that the question they asked cannot be answered by the data they have. The stability-selection and nested-CV instincts are what mark seniority.}

\goodgreat
\begin{itemize}
    \item \badgeLfive{L5} Diagnoses $d>n$ and collinearity, proposes ridge, interprets VIF correctly.
    \item \badgeLsix{L6} Splits the prediction-versus-inference goals, prefers elastic net for correlated groups with the reason, adds stability selection and nested CV.
    \item \badgeLseven{L7} Prices the alternatives---collect more data, reduce the feature space by domain knowledge, or reframe the question as a designed experiment (Chapter~\ref{cml:chap:experimentation_causal})---and states plainly what claims the data can and cannot support.
\end{itemize}

\followups{``What exactly does VIF $=40$ mean?'' ``Why does ridge have a solution when OLS does not?'' ``How would you present feature importance to a stakeholder here?''}
\end{interviewq}

\begin{interviewq}
\textbf{Q: You need to model weekly support tickets per account. The target is a non-negative integer, most accounts have zero or one, a few have hundreds, and accounts have very different sizes. Design the model.}

\badgeLsix{L6} \quad \badgetype{Architecture Design} \quad \badgefreq{\freqlow}

\stronganswer

Start with the response distribution, not the algorithm. The target is a \textbf{count} with a long right tail and unequal exposure, which rules out plain OLS: it can predict negative values, its constant-variance assumption is badly violated (count variance grows with the mean), and it is dominated by the few huge accounts. Two common patches are also wrong---modeling $\log(y+1)$ with OLS estimates $\E[\log y]$, not $\log \E[y]$, and back-transforming needs a smearing correction, while the $+1$ is an arbitrary distortion of the many zeros. \textbf{The right frame is a Poisson GLM with a log link and an offset}: $\log\E[y_i] = \log(\text{exposure}_i) + x_i^{\top}\beta$, where exposure is seats, monthly active users, or account-days. The offset is a coefficient fixed at $1$, which encodes ``twice the size, twice the expected tickets'' and turns the model into one for a \emph{rate}; forcing the model to learn that from a size feature instead wastes capacity and generalizes worse to unseen sizes. Coefficients then read as multiplicative rate ratios, $e^{\beta_j}$, which is a good artifact for stakeholders. \textbf{Then check the Poisson assumption}, because it forces $\Var=\mu$ and support tickets are almost always overdispersed (a single incident produces a burst). Diagnose by comparing residual deviance to residual degrees of freedom, or by binning predictions and plotting empirical variance against mean. If overdispersed: negative binomial (adds a real dispersion parameter and widens the intervals appropriately) or quasi-Poisson (same point estimates, corrected standard errors) when only inference is at stake. \textbf{Check zero inflation separately}: if there are more exact zeros than even a negative binomial predicts---typically because a subset of accounts \emph{never} files tickets---use a hurdle or zero-inflated model, which is honest about two distinct processes. \textbf{Validation and metrics}: evaluate with Poisson deviance or the appropriate log-likelihood, not RMSE, and split by \emph{time} and by account (the same account across weeks is not independent---use grouped/blocked CV or cluster-robust errors). \textbf{Finally, the model-class question}: if the goal is pure predictive accuracy and there are enough rows, a GBDT with a Poisson objective and the exposure as an offset usually beats the GLM (Chapter~\ref{cml:chap:trees_ensembles})---the GLM's advantages are interpretability, small-data stability, and clean uncertainty. Say which you are optimizing for and why.

\redflags
\begin{itemize}
    \item Reaching for linear regression on a count, or for $\log(y+1)$ with OLS as though it were equivalent.
    \item Treating it as binary (``will this account file a ticket?'') without being asked to.
    \item No exposure/offset, so a 10-seat and a 10{,}000-seat account are modeled as comparable.
    \item Never checking overdispersion, then reporting confidence intervals that are far too narrow.
    \item Evaluating with RMSE and being surprised that the model ignores small accounts.
\end{itemize}

\interviewtest{Whether you choose a likelihood from the data-generating process rather than reaching for a default regressor---the practical payoff of the GLM frame. The offset is the specific detail interviewers listen for; overdispersion is the second.}

\goodgreat
\begin{itemize}
    \item \badgeLfive{L5} Poisson regression with a log link, and non-negativity as the reason not to use OLS.
    \item \badgeLsix{L6} Adds the exposure offset with its interpretation, the overdispersion diagnosis and the negative-binomial/quasi-Poisson fix, and grouped temporal validation with deviance as the metric.
    \item \badgeLseven{L7} Adds zero inflation as a distinct hypothesis with its own test, states when to abandon the GLM for a Poisson-objective GBDT and what is lost, and connects the choice to the downstream decision (staffing, which needs a predictive interval, not a point estimate).
\end{itemize}

\followups{``What exactly does the offset do to the likelihood?'' ``How do you detect overdispersion?'' ``When is a hurdle model better than zero-inflated Poisson?''}
\end{interviewq}

\begin{interviewq}
\textbf{Q: Design the model layer for an ad click-through-rate system: $10^9$ events per day, $10^8$ hashed sparse features, hourly retraining, sub-millisecond scoring, and the score is multiplied by a bid.}

\badgeLseven{L7} \quad \badgetype{System Design} \quad \badgefreq{\freqlow}

\stronganswer

Lead with the constraints that pick the model, not with the model. Four bind here: \emph{feature space} ($10^8$ sparse, dominated by ultra-high-cardinality IDs and crosses), \emph{freshness} (hourly or continuous---the distribution moves within a day), \emph{latency} (sub-millisecond, so the serving path must be a sparse dot product), and \emph{calibration} (the score is multiplied by a bid, so a $20\%$ systematic overestimate is a $20\%$ overspend---this system needs \emph{probabilities}, not rankings). \textbf{Model.} Logistic regression on hashed sparse features trained with \textbf{FTRL-Proximal} is the canonical answer and still a live production pattern: one-pass online learning so the model absorbs new data continuously, per-coordinate adaptive learning rates so that rare features get their own effective step size, and L1 in the proximal step producing \emph{genuine} zeros---which matters operationally because the served model's memory footprint is the number of nonzero weights. The feature hashing trick fixes the parameter space at a chosen size and handles unbounded vocabulary growth (new ads, new URLs) at the cost of collisions, which are tolerable at this scale and can be measured. \textbf{Where nonlinearity comes from}: explicit feature crosses, and the classic GBDT-leaf-indices-as-features pattern, both of which keep the serving path linear. When the team is ready to pay for it, the modern answer is a two-tower or Wide-and-Deep-style model with embeddings for the high-cardinality IDs and a linear/sparse wide part for memorization---but that changes the serving story and the retraining cost, and it should be justified by a measured lift, not by novelty [DL 12], [SR 6]. \textbf{Calibration is a first-class component}, not an afterthought: train on negatively downsampled data for tractability and \emph{correct the intercept} analytically for the sampling rate; monitor calibration continuously by slice (ad, placement, device, hour); recalibrate with isotonic regression when drift appears. Report expected calibration error and reliability curves alongside AUC, because AUC is invariant to any monotone rescaling and will look perfect while the bidder loses money [DL 23], Chapter~\ref{cml:chap:probabilistic_foundations}. \textbf{Serving and safety}: sparse weights pinned in memory, feature fetching (not the dot product) as the real p99 driver, shadow scoring and A/B for every model change, automatic rollback on calibration or volume alarms, and a temporal---never random---evaluation split, since a random split leaks the future and every offline number becomes optimistic (Chapter~\ref{cml:chap:applied_craft}, Chapter~\ref{cml:chap:experimentation_causal}).

\redflags
\begin{itemize}
    \item Jumping straight to a deep model without addressing freshness, latency, or calibration---the constraints that actually decide this.
    \item Treating CTR as a ranking problem. The score is multiplied by a bid; miscalibration is a direct revenue error.
    \item Proposing daily batch retraining on a distribution that moves hourly, or a random train/test split on timestamped data.
    \item Downsampling negatives without correcting the intercept, which biases every predicted probability upward.
    \item No monitoring plan---this is a system that fails silently and expensively.
\end{itemize}

\interviewtest{Whether classical-model knowledge survives contact with production constraints: online learning, sparsity as a memory budget, calibration as a revenue property, and temporal validation. The question rewards knowing why the simple model is the right one here.}

\goodgreat
\begin{itemize}
    \item \badgeLfive{L5} Logistic regression on hashed features with online updates, and calibration named.
    \item \badgeLsix{L6} FTRL with per-coordinate rates and L1-for-model-size, the downsampling intercept correction, temporal splits, and slice-level calibration monitoring.
    \item \badgeLseven{L7} Prices the migration to embedding models against its serving and retraining cost, designs the rollout and rollback criteria, and identifies the failure modes the metrics will not catch (position bias in the training clicks, feedback loops from the bidder's own decisions [SR 5]).
\end{itemize}

\followups{``Derive the intercept correction for negative downsampling.'' ``AUC is flat but revenue dropped---what do you check?'' ``When does the embedding model finally pay for itself?''}
\end{interviewq}
